\chapter{Extragalactic Astrophysics}

\section*{Theory problems}

\probhead{12.1}{}{2008}{Bandung, Indonesia}{TH S15}{20}
% source id: 2008_TH_S15   tag: Tully-Fisher mass and luminosity NGC 2639
Galaxy NGC 2639 is morphologically identified as an Sa galaxy with measured
maximum rotational velocity $v_{\max}$ of 324\,km/s. After corrections for any
extinction, its apparent magnitude in $B$ is $m_B = 12.22$. It is customary to
measure a radius $R_{25}$ (in units of kpc) at which the galaxy's surface
brightness falls to $25\,\mathrm{mag}_B/\mathrm{arcsec}^2$. Spiral galaxies
tend to follow a typical relation:
\[ \log R_{25} = -0.249 M_B - 4.00, \]
where $M_B$ is the absolute magnitude in $B$. Apply the $B$-band Tully--Fisher
relation for Sa spirals
\[ M_B = -9.95 \log v_{\max} + 3.15 \quad (v_{\max}\ \text{in km/s}) \]
to calculate the mass of NGC 2639 out to $R_{25}$. If colour index of the sun
is $(m_{B_\odot} - m_{V_\odot}) = 0.64$, write the mass (of NGC 2639) in units
of solar mass $M_\odot$ and its luminosity $B$-band in unit of $L_\odot$.

\probhead{12.2}{}{2011}{Katowice, Poland}{TH L2}{30}
% source id: 2011_TH_L2   tag: elliptical galaxy K-correction, cluster membership
Within the field of a galaxy cluster at a redshift of $z = 0.500$, a galaxy
which looks like a normal elliptical is observed, with an apparent magnitude in
the $B$ filter $m_B = 20.40$\,mag.

The luminosity distance corresponding to a redshift of $z = 0.500$ is
$d_L = 2754$\,Mpc.

The spectral energy distribution (SED) of elliptical galaxies in the wavelength
range 250\,nm to 500\,nm is adequately approximated by the formula:
\[ L_\lambda(\lambda) \propto \lambda^4 \]
(i.e., the spectral density of the object's luminosity, known also as the
monochromatic luminosity, is proportional to $\lambda^4$.)

\begin{description}
\item[(a)] What is the absolute magnitude of this galaxy in the $B$ filter?
\item[(b)] Can it be a member of this cluster? (write YES or NO alongside your
  final calculation)
\end{description}

Hints: Try to establish a relation that describe the dependence of the spectral
density of flux on distance for small wavelength interval. Normal elliptical
galaxies have maximum absolute magnitude equal to $-22$\,mag.

\probhead{12.3}{}{2012}{Rio de Janeiro, Brazil}{TH L2}{}
% source id: 2012_TH_L2   tag: spiral rotation curve, Tully-Fisher
Astronomers studied a spiral galaxy with an inclination angle of $90^\circ$
from the plane of the sky (``edge-on'') and apparent magnitude 8.5. They
measured the rotational velocity and radial distance from the Galactic center
and plotted its rotation curve.

\begin{description}
\item[(2.1)] Approximate the rotation curve in Figure 1 with a continuous
  function $V(D)$ composed of two straight lines.
\item[(2.2)] Using the same observations, they estimated that the rotation
  period of the pressure wave in the galactic disk is half of the rotation
  period of the mass of the disk. Estimate the time it takes for a spiral arm
  to take another turn around the galactic center (use the function
  constructed in 2.1).
\item[(2.3)] Calculate the distance to the galaxy using the Tully--Fisher
  relation (see Table of Constants).
\item[(2.4)] Calculate the maximum and minimum values of the observed
  wavelengths of the hydrogen lines corresponding to 656.28\,nm in the
  spectrum of this galaxy. Hint: also take into account the cosmological
  expansion.
\item[(2.5)] Using Figure 1, estimate the mass of the galaxy up to a radius of
  $3 \times 10^3$ light-years.
\item[(2.6)] Estimate the number of stars of the galaxy, assuming that:
  \begin{itemize}
  \item the mean mass of the stars is equal to one solar mass and one third of
    the baryonic mass of the galaxy is in the form of stars, and;
  \item the fraction of baryonic to dark matter in the galaxy is the same as
    the fraction for the whole Universe (see Table of Constants).
  \end{itemize}
\end{description}
\ioaafig{2012_TH_L2-f3.pdf}

\probhead{12.4}{}{2015}{Magelang, Indonesia}{TH S7}{}
% source id: 2015_TH_S7   tag: galaxy cluster density from escape velocity
A galaxy at the boundary of a galaxy cluster of radius 10\,Mpc is expected to
escape from the cluster if it has an initial velocity of at least 700\,km/s
relative to the center of the cluster. Calculate the density of the cluster.

\probhead{12.5}{Mass of the Local Group}{2017}{Phuket, Thailand}{TH T11}{50}
% source id: 2017_TH_T11   tag: MW-M31 timing argument
The dynamics of M31 (Andromeda) and the Milky Way (MW) can be used to estimate
the total mass of the Local Group (LG). The basic idea is that galaxies
currently in a binary system were at approximately the same point in space
shortly after the Big Bang. To a reasonable approximation, the mass of the
local group is dominated by the masses of the MW and M31. Via Doppler shifts
of the spectral lines, it was found that M31 is moving towards the MW with a
speed of 118\,km\,s$^{-1}$. This may be surprising, given that most galaxies
are moving away from each other with the general Hubble flow. The fact that
M31 is moving towards the MW is presumably because their mutual gravitational
attraction has eventually reversed their initial velocities. In principle, if
the pair of galaxies is well-represented by isolated point masses, their total
mass may be determined by measuring their separation, relative velocity and
the time since the universe began. Kahn and Woltjer (1959) used this argument
to estimate the mass in the LG.

In this problem we will follow this argument through our calculation as
follows.

\begin{description}
\item[(a)] Consider an isolated system with negligible angular momentum of two
  gravitating point masses $m_1$ and $m_2$ (as observed by an inertial
  observer at the centre of mass).

  \ioaafig{2017_TH_T11-f1.pdf}

  Write down the expression of the total mechanical energy ($E$) of this
  system in mathematical form connecting $m_1$, $m_2$, $r_1$, $r_2$, $v_1$,
  $v_2$, and the universal gravitational constant $G$, where $v_1$ and $v_2$
  are the radial velocities of $m_1$ and $m_2$, respectively. \hfill[5]
\item[(b)] Re-write the equation in a) in terms of $r$, $v$, $\mu$, $M$, and
  $G$, where $r \equiv r_1 + r_2$ is the separation distance between $m_1$ and
  $m_2$, $v$ is the changing rate of the separation distance,
  $\mu \equiv \dfrac{m_1 m_2}{m_1 + m_2}$ is the reduced mass of the system,
  and $M \equiv m_1 + m_2$ is the total mass of the system. \hfill[10]
\item[(c)] Show that the equation in b) yields
  \[ v^2 = (2GM)\left(\frac{1}{r} - \frac{1}{r_0}\right),
     \quad\text{where } r_0 \text{ is a new constant.} \]
  Find $r_0$ in terms of $\mu$, $M$, $G$ and $E$. \hfill[5]
\end{description}

The solution of the equation in b) is given below in parametric form, under
the initial condition $r = 0$ at $t = 0$:
\[ r(\theta) = \frac{r_0}{2}\left(1 - \cos\theta\right), \]
\[ t(\theta) = \left(\frac{r_0^3}{8GM}\right)^{\frac{1}{2}}
   \left(\theta - \sin\theta\right), \]
where $\theta$ is in radians.

\begin{description}
\item[(d)] From the above parametric equations, show that an expression for
  $\dfrac{vt}{r}$ is
  \[ \frac{vt}{r} = \frac{(\sin\theta)(\theta - \sin\theta)}
     {(1 - \cos\theta)^2} \]
  \hfill[10]
\item[(e)] Now we consider $m_1$ and $m_2$ as the MW and M31, respectively,
  such that the current values of $v$ and $r$ are $v = -118$\,km\,s$^{-1}$ and
  $r = 710$\,kpc, and $t$ may be taken to be the age of the Universe (13700
  million years). Find $\theta$ using numerical iteration. \hfill[10]
\item[(f)] Use the value of $\theta$ from e) to calculate the maximum distance
  between M31 and the MW, $r_{\max}$, and hence also obtain the value of $M$
  in solar masses. \hfill[10]
\end{description}

\probhead{12.6}{Galactic Outflow}{2017}{Phuket, Thailand}{TH T9}{20}
% source id: 2017_TH_T9   tag: galaxy rotation, escape velocity, outflow
Cannon et al.\ (2004) conducted an HI observation of a disk starburst galaxy,
IRAS 0833+6517, with the Very Large Array (VLA). The galaxy is located at a
distance of 80.2\,Mpc with an approximate inclination angle of 23 degrees.
According to the HI velocity map, IRAS 0833+6517 appears to be undergoing
regular rotation with an observed radial velocity of the HI gas of roughly
5850\,km\,s$^{-1}$ at a distance of 7.8\,kpc from the centre (the left panel
of the figure below).

Gas outflow from IRAS 0833+6517 is traced by using the blueshifted
interstellar absorption lines observed against the backlight of the stellar
continuum (the right panel of the figure). Assuming that this galaxy is
gravitationally stable and all the stars are moving in circular orbits,

\begin{description}
\item[(a)] Determine the rotational velocity ($v_{\mathrm{rot}}$) of IRAS
  0833+6517 at the observed radius of HI gas. \hfill[5]
\item[(b)] Calculate the escape velocity for a test particle in the gas
  outflow at the radius of 7.8\,kpc. \hfill[9]
\item[(c)] Determine if the outflowing gas can escape from the galaxy at this
  radius by considering the velocity offset of the C\,II $\lambda$1335
  absorption line, which is already corrected for the cosmological recessional
  velocity. (The central rest-frame wavelength of the CII absorption line is
  1335\,\AA.) (YES / NO) \hfill[6]
\end{description}
\ioaafig{2017_TH_T9-f1.pdf}

\probhead{12.7}{A Possible Dark Matter Deficient Galaxy}{2018}{Beijing, China}{TH T8}{25}
% source id: 2018_TH_T8   tag: NGC1052-DF2 dynamical vs stellar mass
Earlier this year, a team of astronomers reported their discovery of a galaxy
with much less dark matter than the galaxy evolution model predicted (van
Dokkum et al.\ 2018, Nature). This galaxy, named NGC 1052-DF2, is located
close to the elliptical galaxy NGC 1052 ($D = 20$\,Mpc from the Sun) in the
sky. The shape of NGC 1052-DF2 resembles an ellipse with semi major axis ($a$)
of $22.6''$ and $\frac{b}{a} = 0.85$. Half of the total light from the galaxy
comes from within this ellipse and the mean surface brightness within the
ellipse is about 24.7\,mag\,arcsec$^{-2}$.

\begin{description}
\item[(a)] Calculate the total apparent magnitude of this galaxy.
\item[(b)] The team suggested the galaxy is a companion of NGC 1052.
  Determine the total mass of stars in NGC 1052-DF2, assuming it has a mass to
  light ratio $\left(\frac{M/M_\odot}{L/L_\odot}\right)$ of 2.0.
\item[(c)] The team identified 10 globular clusters in NGC 1052-DF2 with a
  mean galactocentric distance of $78.4''$. They also measured the velocity
  dispersion of these clusters to be not more than 8.4\,km/s. Estimate the
  dynamical mass of this galaxy. For simplicity, assume the mass distribution
  in the galaxies is uniform and is spherically symmetric.
\item[(d)] This discovery was challenged by other groups (Kroupa et al.,
  Nature, 2018, Truijlo et al., MNRAS, 2018), who claimed that NGC 1052-DF2 is
  not a satellite of NGC 1052, and it is located at a much smaller distance to
  us. Show why a smaller distance would weaken the assertion of the dark
  matter deficiency in NGC 1052-DF2.
\end{description}

\probhead{12.8}{X-ray emission from galaxy clusters}{2023}{Katowice, Poland}{TH Q11}{30}
% source id: 2023_TH_Q11   tag: bremsstrahlung plasma mass, Sunyaev-Zeldovich
Clusters of galaxies are strong X-ray sources. It has been established that
the emission mechanism is thermal bremsstrahlung (free-free radiation) from a
hot hydrogen and helium plasma inside the cluster. The luminosity $L_X$ (in
Watts) of each component of the plasma is described by the formula:
\[ L_X = 6 \times 10^{-41}\, N_e\, N_X\, T^{\frac{1}{2}}\, V Z_X^2, \]
where the symbols represent:
\begin{itemize}
\item $X$ -- Hydrogen (H) or Helium (He),
\item $N_e$ -- number density of electrons [$\mathrm{m}^{-3}$],
\item $N_X$ -- number density of ions $X$ [$\mathrm{m}^{-3}$],
\item $Z_X$ -- atomic number of ion $X$,
\item $T$ -- temperature of the plasma [K],
\item $V$ -- volume occupied by the plasma [$\mathrm{m}^3$].
\end{itemize}

\begin{description}
\item[(a)] Determine the total mass (in solar masses) of the plasma which
  emits the X-rays, assuming that:
  \begin{itemize}
  \item the plasma is fully ionized with 1 helium ion for every 10 hydrogen
    ions,
  \item $L_{\mathrm{total}} = 1.0 \times 10^{37}$\,W,
  \item $T = 80 \times 10^6$\,K,
  \item the plasma is uniformly distributed in a sphere of radius
    $R = 500$\,kpc,
  \item self-absorption is negligible.
  \end{itemize}
  \hfill(16 points)
\end{description}

The photons of the cosmic microwave background (CMB) interact with plasma in a
process known as inverse Compton scattering. The CMB normally has a thermal
blackbody spectrum at a temperature of 2.73\,K. However, interaction with the
plasma leads to distortion of the CMB spectrum (known as the
Sunyaev-Zeldovich effect).

\begin{description}
\item[(b)] Estimate the mean free path of CMB photons in the plasma, i.e.\ the
  average distance travelled by a photon before interacting with an electron.
  Express it in Mpc. The effective cross section for photon-electron
  interactions is $\sigma = 6.65 \times 10^{-29}$\,$\mathrm{m}^2$.
  \hfill(5 points)
\item[(c)] Estimate the typical energy of CMB photons. \hfill(3 points)
\item[(d)] The energy of CMB photons can be increased by a factor of up to
  $(1+\beta)/(1-\beta)$ due to the inverse Compton scattering, where
  $v = \beta c$ is the velocity of electrons. Estimate the energy of scattered
  CMB photons. \hfill(6 points)
\end{description}

\probhead{12.9}{Galaxy Cluster}{2024}{Rio de Janeiro, Brazil}{TH T2}{10}
% source id: 2024_TH_T2   tag: virial cluster mass from redshift dispersion
An astrophysical survey mapped all the galaxies in a small region of the sky,
of angular diameter $\Delta\theta = 0.01$\,rad, where many galaxies seemed to
be concentrated around the central area of the image. When the positions and
redshifts of all the galaxies in this cluster were measured, an interesting
distribution emerged, which is shown in the plot below.

\ioaafig{2024_TH_T2-f1.pdf}

Using these observations, estimate the total mass of the galaxy cluster and
express your answer in solar masses. Assume that this galaxy cluster is in
dynamical equilibrium, with a root-mean-square redshift dispersion
$\sigma_z = \sqrt{\langle (z - 0.7)^2 \rangle} = 0.0005$. Feel free to make
reasonable approximations when considering the average velocities, masses, and
spatial distribution of the galaxies.

Consider that the distance to $\bar{z} = 0.7$ in the standard cosmological
model is $D_A = 1500$\,Mpc. Ignore cosmological effects on the distance.

\probhead{12.10}{Quasars}{2025}{Mumbai, India}{TH T05}{20}
% source id: 2025_TH_T05   tag: superluminal jet motion, core size
A quasar is an extremely luminous active galaxy powered by a supermassive
black hole that emits relativistic jets. The figure shows a series of panels
of radio images of a quasar (with redshift $z = 0.53$, and luminosity distance
$D_{\mathrm{L}} = 1.00 \times 10^{10}$\,ly) at different times. The ``core''
aligns with the vertical white line, while a jet, consisting of a ``blob''
(marked white $+$), moves away from it over time. Each panel shows the
observation time (starting with $T_0$ for the first image), and the angular
scale is indicated at the top and bottom of the figure.

\ioaafig{2025_TH_T05-f1-panel1.pdf}
\ioaafig{2025_TH_T05-f1-panel2.pdf}
\ioaafig{2025_TH_T05-f1-panel3.pdf}
\ioaafig{2025_TH_T05-f1-panel4.pdf}
\ioaafig{2025_TH_T05-f1-panel5.pdf}

\begin{description}
\item[(T05.1)] Determine the blob's angular separation, $\phi_{\mathrm{blob}}$
  (in milliarcsecond), and its transverse distance, $l_{\mathrm{blob}}$ (in
  light-year), from the quasar core for each observation. Then, calculate the
  blob's apparent velocity in the transverse direction ($v_{\mathrm{app}}$) as
  a fraction of light speed, $\beta_{\mathrm{app}}$
  ($= v_{\mathrm{app}}/c$) by using consecutive observations. Also calculate
  the average apparent velocity $\beta_{\mathrm{app}}^{\mathrm{ave}}$ over the
  entire observation period. \hfill[5]
\end{description}

The quasar jet actually moves at a relativistic speed $v \equiv \beta c$, but
not necessarily in the plane of the sky; e.g., it makes an angle $\theta$ (the
``viewing angle'') with respect to the line of sight of a distant observer
(indicated by the dashed lines), as shown in the sketch below.

For this and all subsequent parts, ignore redshift of the quasar and any
relativistic effects.

\ioaafig{2025_TH_T05-f2.pdf}

\begin{description}
\item[(T05.2)] The light emitted by the blob at two different times $t_0$
  (corresponding to position A) and $t_0 + \Delta t$ (corresponding to
  position B) reaches the observer at $t_{\mathrm{A}}$ and $t_{\mathrm{B}}$,
  respectively. Thus the observed time difference is
  $\Delta t_{\mathrm{app}} = t_{\mathrm{B}} - t_{\mathrm{A}}$.
  \begin{description}
  \item[(T05.2a)] Find an expression for the ratio
    $\dfrac{\Delta t_{\mathrm{app}}}{\Delta t}$ in terms of $\beta$ and
    $\theta$. \hfill[2]
  \item[(T05.2b)] Using this ratio, express $\beta_{\mathrm{app}}$ in terms of
    $\beta$ and $\theta$. \hfill[2]
  \end{description}
\item[(T05.3)] Motion is called superluminal if the apparent speed exceeds
  that of light ($\beta_{\mathrm{app}} > 1$), and subluminal if it does not
  ($\beta_{\mathrm{app}} < 1$).
  \begin{description}
  \item[(T05.3a)] For $\beta_{\mathrm{app}} = 1$, plot a smooth curve of
    $\beta$ as a function of $\theta$ to mark the boundary between subluminal
    and superluminal motions. Shade the superluminal region in the graph with
    slanted lines. \hfill[4]
  \item[(T05.3b)] Find the lowest true jet speed
    ($\beta_{\mathrm{low}} = v_{\mathrm{low}}/c$) for the superluminal motion
    to occur and also its corresponding viewing angle $\theta_{\mathrm{low}}$.
    \hfill[2]
  \end{description}
\item[(T05.4)] Find an expression for the maximum viewing angle,
  $\theta_{\max}$, for which a given value of $\beta_{\mathrm{app}}$ will be
  possible. \hfill[2]
\end{description}

The core of a quasar, its central compact object, exhibits variability in its
emission due to internal processes occurring within a causally connected
region. The size ($=$ radius) of this region is typically taken to be about
five times the Schwarzschild radius of the core.

\begin{description}
\item[(T05.5)] The core of a certain quasar is found to vary on time scales of
  about 1\,h. Obtain an upper limit, $M_{\mathrm{c,max}}$, on the mass of the
  central compact object, in units of solar mass. \hfill[3]
\end{description}

\section*{Data-analysis problems}

\probhead{12.11}{Weighing a galaxy}{2011}{Katowice, Poland}{DA D2}{25}
% source id: 2011_DA_D2   tag: NGC 7083 rotation curve, bulge mass
The attached images show a photograph of the spiral galaxy NGC 7083, which
lies at a distance of 40\,Mpc, and a fragment of its spectrum. The slit of the
spectrograph was aligned with the major axis of the image of the galaxy. The
$x$-axis of the spectrum represents wavelength, and the $y$-axis represents
the angular distance of the emitting region from the core of the galaxy, where
1\,pixel $= 0.82$\,arcsec. Two bright emission lines are visible, with rest
wavelengths of $\lambda_1 = 6564$\,\AA, $\lambda_2 = 6584$\,\AA.

Use the spectrum to plot the rotation curve of the galaxy and estimate the
mass of the central bulge.

Assumption: central bulge is spherical.

The photograph of the galaxy has the correct proportions.
\ioaafig{2011_DA_D2-f1.pdf}
\ioaafig{2011_DA_D2-f2.pdf}

\probhead{12.12}{Type IA Supernovae}{2016}{Bhubaneswar, India}{DA D3}{50}
% source id: 2016_DA_D3   tag: SN Ia standardization, H0, universe age
Supernovae of type Ia are considered very important for the measurements of
large extragalactic distances. The brightening and subsequent dimming of these
explosions follow a characteristic light curve, which helps in identifying
these as supernovae of type Ia.

Light curves of all type Ia supernovae can be fit to the same model light
curve, when they are scaled appropriately. In order to achieve this, we first
have to express the light curves in the reference frame of the host galaxy by
taking care of the cosmological stretching/dilation of all observed time
intervals, $\Delta t_{\mathrm{obs}}$, by a factor of $(1+z)$. The time
interval in the rest frame of the host galaxy is denoted by
$\Delta t_{\mathrm{gal}}$.

The rest frame light curve of a supernova changes by two magnitudes compared
to the peak in a time interval $\Delta t_0$ after the peak. If we further
scale the time intervals by a factor of $s$ (i.e.\
$\Delta t_s = s\,\Delta t_{\mathrm{gal}}$) such that the scaled value of
$\Delta t_0$ is the same for all supernovae, the light curves turn out to have
the same shape. It also turns out that $s$ is related linearly to the absolute
magnitude, $M_{\mathrm{peak}}$, at the peak luminosity for the supernova. That
is, we can write
\[ s = a + b M_{\mathrm{peak}}, \]
where $a$ and $b$ are constants. Knowing the scaling factor, one can determine
absolute magnitudes of supernovae at unknown distances from the above linear
equation.

The table below contains data for three supernovae, including their distance
moduli, $\mu$ (for the first two), their recession speed, $cz$, and their
apparent magnitudes, $m_{\mathrm{obs}}$, at different times. The time
$\Delta t_{\mathrm{obs}} \equiv t - t_{\mathrm{peak}}$ shows number of days
from the date at which the respective supernova reached peak brightness. The
observed magnitudes have already been corrected for interstellar as well as
atmospheric extinction.

\begin{center}
\begin{tabular}{c|c|c|c}
Name & SN2006TD & SN2006IS & SN2005LZ \\
$\mu$ (mag) & 34.27 & 35.64 & \\
$cz$ (km\,s$^{-1}$) & 4515 & 9426 & 12060 \\
\hline
$\Delta t_{\mathrm{obs}}$ (days) & $m_{\mathrm{obs}}$ (mag)
  & $m_{\mathrm{obs}}$ (mag) & $m_{\mathrm{obs}}$ (mag) \\
\hline
$-15.00$ & 19.41 & 18.35 & 20.18 \\
$-10.00$ & 17.48 & 17.26 & 18.79 \\
$-5.00$  & 16.12 & 16.42 & 17.85 \\
0.00     & 15.74 & 16.17 & 17.58 \\
5.00     & 16.06 & 16.41 & 17.72 \\
10.00    & 16.72 & 16.82 & 18.24 \\
15.00    & 17.53 & 17.37 & 18.98 \\
20.00    & 18.08 & 17.91 & 19.62 \\
25.00    & 18.43 & 18.39 & 20.16 \\
30.00    & 18.64 & 18.73 & 20.48 \\
\end{tabular}
\end{center}

\begin{description}
\item[(D3.1)] Compute $\Delta t_{\mathrm{gal}}$ values for all three
  supernovae, and fill them in the given blank boxes in the data tables on the
  BACK side of the Summary Answersheet. On a graph paper, plot the points and
  draw the three light curves in the rest frame (mark your graph as ``D3.1'').
  \hfill[15]
\item[(D3.2)] Take the scaling factor, $s_2$, for the supernova SN2006IS to be
  1.00. Calculate the scaling factors, $s_1$ and $s_3$, for the other two
  supernovae SN2006TD and SN2005LZ, respectively, by calculating $\Delta t_0$
  for them. \hfill[5]
\item[(D3.3)] Compute the scaled time differences, $\Delta t_s$, for all three
  supernovae. Write the values for $\Delta t_s$ in the same data tables on the
  Summary Answersheet. On another graph paper, plot the points and draw the 3
  light curves to verify that they now have an identical profile (mark your
  graph as ``D3.3''). \hfill[14]
\item[(D3.4)] Calculate the absolute magnitudes at peak brightness,
  $M_{\mathrm{peak},1}$, for SN2006TD and $M_{\mathrm{peak},2}$, for SN2006IS.
  Use these values to calculate $a$ and $b$. \hfill[6]
\item[(D3.5)] Calculate the absolute magnitude at peak brightness,
  $M_{\mathrm{peak},3}$, and distance modulus, $\mu_3$, for SN2005LZ.
  \hfill[4]
\item[(D3.6)] Use the distance modulus $\mu_3$ to estimate the value of
  Hubble's constant, $H_0$. Further, estimate the characteristic age of the
  universe, $T_{\mathrm{H}}$. \hfill[6]
\end{description}
\ioaafig{2016_DA_D3-f1.pdf}

\probhead{12.13}{Calibrating the distance ladder to the LMC}{2017}{Phuket, Thailand}{DA D1}{75}
% source id: 2017_DA_D1   tag: Cepheid PL/Wesenheit, LMC distance
An accurate trigonometric parallax calibration for Galactic Cepheids has long
been sought, but is very difficult to achieve in practice. All known classical
(Galactic) Cepheids are more than 250\,pc away, therefore for direct distance
estimates to achieve an uncertainty of up to 10\%, parallax uncertainties of
up to $\pm 0.2$\,milliarcsec are needed, requiring space-based observations.
The Hipparcos satellite reported parallaxes for 200 of the nearest Cepheids,
but even the best of these had high uncertainties. Recent progress has come
with the use of the Fine Guidance Sensor on HST where parallaxes (in many
cases) accurate to better than $\pm 10$\% were obtained for 10 Cepheids,
spanning a range of periods from 3.7 to 35.6 days. These nearby Cepheids cover
distances from about 300 to 560\,pc.

The measured periods, $P$, and average magnitudes in $V$, $K$ and $I$ bands
are given in Table~1 as well as the $A_V$ and $A_K$ for extinction in $V$ and
$K$ bands, respectively. The measured parallaxes with their uncertainties are
also given in milliarcsec (mas). All measured apparent magnitudes have
negligible uncertainty.

Table 1: Periods and average apparent magnitudes of 5 Galactic Cepheids with
accurate parallax measurements.

\begin{center}
\begin{tabular}{l|c|c|c|c|c|c|c|c}
 & $P$ & $\langle V\rangle$ & $\langle K\rangle$ & $A_V$ & $A_K$
 & $\langle I\rangle$ & parallax & error \\
 & (day) & (mag) & (mag) & (mag) & (mag) & (mag) & (mas) & (mas) \\
\hline
RT Aur      & 3.728  & 5.464 & 3.925 & 0.20 & 0.02 & 4.778 & 2.40 & 0.19 \\
FF Aql      & 4.471  & 5.372 & 3.465 & 0.64 & 0.08 & 4.510 & 2.81 & 0.18 \\
X Sgr       & 7.013  & 4.556 & 2.557 & 0.58 & 0.07 & 3.661 & 3.00 & 0.18 \\
$\zeta$ Gem & 10.151 & 3.911 & 2.097 & 0.06 & 0.01 & 3.085 & 2.78 & 0.18 \\
l Car       & 35.551 & 3.732 & 1.071 & 0.52 & 0.06 & 2.557 & 2.01 & 0.20 \\
\end{tabular}
\end{center}

\begin{description}
\item[(D1.1)] The observed correlation between the period of a Cepheid and its
  brightness is usually described by the so-called ``Period-Luminosity (PL)
  relation'', where $L \propto P^\beta$. In fact, such a relation is normally
  expressed in terms of the period and absolute magnitude, instead of
  luminosity. Hereafter, we shall refer to the Period-Absolute magnitude
  relation as the conventionally named ``PL relation''.

  Use the data given in Table~1 to plot a suitable linear graph in order to
  derive the Cepheid PL relation for the $V$-band and $K$-band. You should
  plot each graph separately on different pieces of graph paper. Determine the
  slope of the line that best describes the linear relation of the data. (You
  may find the relation
  $\Delta(\log_{10} x) \approx \dfrac{\Delta x}{x \log_e 10}$ useful)
  \hfill[36.5 Marks]
\end{description}

Any apparent differences in PL relations of stars in the different bands can
be explained if one also considers differences in colour. Therefore, the PL
relation is in fact a PLC (Period-Luminosity-Colour) relation. This is from
the reddening effect, due to extinction being a function of wavelength, which
can also vary among different Cepheids due to their different metallicities,
foreground Interstellar Medium and dust.

A new reddening-free magnitude (or bandpass) called ``Wesenheit'' has been
proposed that does not require the explicit information of the extinction of
individual stars but uses colour information from the star itself to get rid
of the effect. For example, $W_{VI}$ use $V$ and $I$ band photometry and is
defined as
\[ W_{VI} = V - \left[\frac{A_V}{E(V-I)}\right](V-I), \]
\[ \phantom{W_{VI}} = V - R_v (V-I) \]
where $R_v$ depends on the reddening law. In this case, we shall take the
value of $R_v$ to be 2.45.

\begin{description}
\item[(D1.2)] From the data given in Table~1, plot and derive the
  reddening-free PL relation using Wesenheit $W_{VI}$ magnitudes. Estimate the
  linear slope of the relation as well as its uncertainty.
  \hfill[14.5 Marks]
\item[(D1.3)] Next, we would like to use the newly-derived PL relations from
  question (D1.1) \& (D1.2) to estimate the distance to the Large Magellenic
  Cloud (LMC) using periods and magnitudes of classical Cepheids in the LMC.
  In Table~2, the periods, average extinction-corrected apparent magnitudes,
  $\langle V_{\mathrm{corr}}\rangle$, and Wesenheit $W_{VI}$ magnitudes are
  given.

  Estimate the distance modulus, $\mu$, to each star and then use all the
  information to derive the distance to the LMC (in parsecs) and its standard
  deviation for each band.

  Compare if the derived distances are statistically different for the 2 bands
  (YES/NO).

  Are the standard deviations of the estimated distances for 2 bands different
  (YES/NO)?

  Based on this dataset, which band ($V$ or Wesenheit) is more accurate in
  estimating the distance to the LMC? \hfill[24 Marks]
\end{description}

Table 2: Period, average extinction-corrected apparent magnitude,
$\langle V_{\mathrm{corr}}\rangle$, and average Wesenheit magnitude
measurements of Cepheids in the LMC.

\begin{center}
\begin{tabular}{l|c|c|c}
 & $P$ & $\langle V_{\mathrm{corr}}\rangle$ & $\langle W_{VI}\rangle$ \\
 & (day) & mag & mag \\
\hline
HV12199 & 2.63  & 16.08 & 14.56 \\
HV12203 & 2.95  & 15.93 & 14.40 \\
HV12816 & 9.10  & 14.30 & 12.80 \\
HV899   & 30.90 & 13.07 & 10.97 \\
HV2257  & 39.36 & 12.86 & 10.54 \\
\end{tabular}
\end{center}
\ioaafig{2017_DA_D1-f1.pdf}
\ioaafig{2017_DA_D1-f2.pdf}

\probhead{12.14}{Dust and Young Stars in Star-forming Galaxies}{2018}{Beijing, China}{DA D1}{75}
% source id: 2018_DA_D1   tag: UV slope, IRX-beta dust attenuation, CR7
As a by-product of the star-forming process in a galaxy, interstellar dust can
significantly absorb stellar light in ultraviolet (UV) and optical bands, and
then re-emit in far-infrared (FIR), which corresponds to a wavelength range of
10--300\,$\mu$m.

1.1. In the UV spectrum of a galaxy, the major contribution is from the light
of the young stellar population generated in recent star-formation processes,
thus the UV luminosity can act as a reliable tracer of the star-formation rate
(SFR) of a galaxy. Since the observed UV luminosity is strongly affected by
dust attenuation, extragalactic astronomers define an index called the
\emph{UV continuum slope} ($\beta$) to quantify the shape of the UV continuum:
\[ f_\lambda = Q \cdot \lambda^\beta \]
where $f_\lambda$ is the monochromatic flux of the galaxy at a given
wavelength $\lambda$ (in the unit of W\,m$^{-3}$) and $Q$ is a scaling
constant.

\begin{description}
\item[(D1.1.1)] (6 points) AB magnitude is a specific magnitude system. The AB
  magnitude is defined as:
  \[ m_{\mathrm{AB}} = -2.5 \log \frac{f_\nu}{3631\,\mathrm{Jy}} \]
  The AB magnitude of a typical galaxy is roughly constant in the UV band.
  What is the UV continuum slope of this kind of galaxy? (Hint:
  $f_\nu \Delta\nu = f_\lambda \Delta\lambda$)
\item[(D1.1.2)] (12 points) Table 1 presents the observed IR photometry
  results for a $z = 6.60$ galaxy called CR7. Plot the AB magnitude of CR7
  versus the logarithm of the rest-frame wavelength on graph paper and
  labelled as Figure 1.
\item[(D1.1.3)] (5 points) Calculate CR7's UV slope, plot the best-fit UV
  continuum on Figure 1 and make a comparison with the results you obtained in
  (D1.1.1). Is it dustier than the typical galaxy in (D1.1.1)? Please answer
  with [YES] or [NO]. (Hint: Express $m_{\mathrm{AB}}$ as a function of
  $\lambda$ and $m_{1600}$, where $m_{1600}$ is the AB magnitude at
  $\lambda_0 = 160$\,nm (1600\,\AA))
\end{description}

Table 1. (Observed Frame) IR Photometry of CR7 at $z = 6.60$.

\begin{center}
\begin{tabular}{l|c|c|c|c}
Band & $Y$ & $J$ & $H$ & $K$ \\
\hline
Central Wavelength ($\mu$m) & 1.05 & 1.25 & 1.65 & 2.15 \\
AB Magnitude & $24.71 \pm 0.11$ & $24.63 \pm 0.13$ & $25.08 \pm 0.14$
  & $25.15 \pm 0.15$ \\
\end{tabular}
\end{center}

1.2. Under the assumption that dust grains in the galaxy absorb the energy of
UV photons and re-emit it by blackbody radiation, the relation between the UV
continuum slope ($\beta$), UV brightness (at 1600\,\AA) and FIR brightness
could be established:
\[ \mathrm{IRX} \equiv \log\!\left(\frac{F_{FIR}}{F_{1600}}\right)
   = S(\beta) \]
where $F_{FIR}$ is the observed far-infrared flux and $F_{1600}$ is the
observed flux at rest-frame wavelength 160\,nm (1600\,\AA) (The ``flux''
$F_\lambda$ is defined as $F_\lambda = \lambda \cdot f_\lambda$). Table 2
presents 20 measurements of $\beta$, $F_{FIR}$ and $F_{1600}$ in nearby
galaxies (Meurer et al.\ 1999).

Table 2. UV slope, flux and FIR flux of 20 nearby galaxies.

\begin{center}
\begin{tabular}{l|c|c|c}
Galaxy Name & UV Slope $\beta$
  & $\log(F_{1600}/10^{-3}\,\mathrm{W\,m}^{-2})$
  & $\log(F_{FIR}/10^{-3}\,\mathrm{W\,m}^{-2})$ \\
\hline
NGC4861      & $-2.46$ & $-9.89$  & $-9.97$  \\
Mrk 153      & $-2.41$ & $-10.37$ & $-10.92$ \\
Tol 1924-416 & $-2.12$ & $-10.05$ & $-10.17$ \\
UGC 9560     & $-2.02$ & $-10.38$ & $-10.41$ \\
NGC 3991     & $-1.91$ & $-10.14$ & $-9.80$  \\
Mrk 357      & $-1.80$ & $-10.58$ & $-10.37$ \\
Mrk 36       & $-1.72$ & $-10.68$ & $-10.94$ \\
NGC 4670     & $-1.65$ & $-10.02$ & $-9.85$  \\
NGC 3125     & $-1.49$ & $-10.19$ & $-9.64$  \\
UGC 3838     & $-1.41$ & $-10.81$ & $-10.55$ \\
NGC 7250     & $-1.33$ & $-10.23$ & $-9.77$  \\
NGC 7714     & $-1.23$ & $-10.16$ & $-9.32$  \\
NGC 3049     & $-1.14$ & $-10.69$ & $-9.84$  \\
NGC 3310     & $-1.05$ & $-9.84$  & $-8.83$  \\
NGC 2782     & $-0.90$ & $-10.50$ & $-9.33$  \\
NGC 1614     & $-0.76$ & $-10.91$ & $-8.84$  \\
NGC 6052     & $-0.72$ & $-10.62$ & $-9.48$  \\
NGC 3504     & $-0.56$ & $-10.41$ & $-8.96$  \\
NGC 4194     & $-0.26$ & $-10.62$ & $-8.99$  \\
NGC 3256     & 0.16    & $-10.32$ & $-8.44$  \\
\end{tabular}
\end{center}

\begin{description}
\item[(D1.2.1)] (14 points) Based on the data given in Table 2, plot the
  $\mathrm{IRX}-\beta$ diagram on graph paper and labelled as Figure 2 and
  find a linear fit to the data. Write down your best-fit equation (i.e.\
  $\mathrm{IRX} = a \cdot \beta + b$) by the side of your diagram.
\item[(D1.2.2)] (6 points) Quantify the dispersion (in `units' of dex, where
  for example, $\log(10^9) - \log(10^4) = 5$\,dex) between the observed
  $\mathrm{IRX}_{\mathrm{obs}}$ and predicted $\mathrm{IRX}_{\mathrm{pred}}$
  using the following equation:
  \[ \sigma = \sqrt{\frac{\sum (\Delta\mathrm{IRX}_i)^2}{N-1}}
     \quad (\text{unit: dex}) \quad \text{where }
     \Delta\mathrm{IRX}_i = \mathrm{IRX}_{i,\mathrm{obs}}
       - \mathrm{IRX}_{i,\mathrm{pred}} \]
\end{description}

1.3. Under the previous assumption of the energy transfer process, the ratio
of $F_{FIR}$ to $F_{1600}$ can be expressed as:
\[ \frac{F_{FIR}}{F_{1600}} \approx 10^{0.4 A_{1600}} - 1 \]
where $F_{1600}$ is the unattenuated flux and $A_\lambda$ is the dust
absorption (in magnitudes) as a function of wavelength $\lambda$.

\begin{description}
\item[(D1.3.1)] (6 points) Express $A_{1600}$ as a function of IRX.
\item[(D1.3.2)] (12 points) Based on Table 2 data and the function of
  $A_{1600}(\mathrm{IRX})$ you derived above, plot the $A_{1600}-\beta$
  diagram on graph paper and label it as Figure 3 and find a linear fit to the
  data. Write down your best-fit equation (i.e.\
  $A_{1600} = a' \cdot \beta + b'$) by the side of your diagram.
\item[(D1.3.3)] (2 points) If your linear model in (D1.3.2) is correct, what
  is the expected UV continuum slope $\beta_0$ of a dust-free galaxy?
\end{description}

1.4. After establishing the local relation between UV continuum slope and IRX,
we could probably test this empirical law in the high-redshift universe. In
2016, researchers obtained an Atacama Large Millimeter\,/\,submillimeter Array
(ALMA) observation of CR7, and the FIR continuum corresponded to a $3\sigma$
upper limit of an FIR flux of $1.5 \times 10^{-19}$\,W\,m$^{-2}$.

\begin{description}
\item[(D1.4.1)] (6 points) Calculate the IRX of CR7. Is it an upper limit or
  lower limit?

  Hint: here $F_{1600}$ should be written in the form of:
  \[ F_{1600} = \lambda_0 \cdot f_{1600} \]
  where $\lambda_0 = 160$\,nm (1600\,\AA) and $f_{1600}$ is the observed flux
  in the rest-frame.
\item[(D1.4.2)] (6 points) Is the current observation long enough to show any
  deviation of CR7 from the $\mathrm{IRX}-\beta$ relationship you just derived
  in the local universe? Please answer with [YES] or [NO] on the summary
  answer sheet, give the IRX difference and show the working used to calculate
  it on the answer sheet.
\end{description}
\ioaafig{2018_DA_D1-f1.pdf}

\probhead{12.15}{AGN}{2020}{GeCAA (Estonia, remote)}{DA 1}{15}
% source id: 2020_DA_1   tag: reverberation mapping, virial BH mass
It is believed that the accretion disk around supermassive black holes (BH) at
galactic centres gives rise to UV thermal emission. This emission is
associated with Active Galactic Nuclei (AGNs).

The optical spectra of bright AGNs show additional bright broad emission
lines. Those emission lines arise from the dense gas in the Broad Line Region
(BLR), which is ionized by the UV photons from the accretion disk. See the
sketch to visualise this model.

\ioaafig{2020_DA_1-f1.pdf}

We can assume that the flux of broad emission lines varies in response to the
variation of the UV continuum with a time delay. This time delay should be
proportional to the separation $R_{\mathrm{BLR}}$ between the BH and the BLR.

Assume that the size of the accretion disk is negligible as compared to
$R_{\mathrm{BLR}}$.

\begin{description}
\item[(a)] (1 point) Estimate the time lag (days) between the B-band continuum
  and broad emission line (H-$\beta$) using the light curves shown below. The
  x-axis is in reduced Julian Dates (JD).

  \ioaafig{2020_DA_1-f2.pdf}
\item[(b)] (3 points) Estimate $R_{\mathrm{BLR}}$ in parsecs (pc).
\item[(c)] (2 points) Estimate the angular separation of this region
  $\theta_{\mathrm{BLR}}$ (in arcsec) from the blackhole, if this AGN is
  100\,Mpc away from us.

  It is possible to estimate the mass of the system using the Virial theorem,
  if the velocity dispersion of the gasses in the BLR and the size of the
  system are known. Assume that the masses of the accretion disk and broad
  line region are negligible, as compared to the black hole. The velocity
  dispersion $v_\sigma$ may be estimated from the broadening of the given
  emission line. We will take the corresponding wavelength dispersion to be
  \[ \sigma = \frac{\mathrm{FWHM}}{2.35} \]
  where FWHM is the full width at half maximum of the broad emission line.
\item[(d)] (5 points) Calculate the velocity dispersion $v_\sigma$ in units of
  km\,s$^{-1}$, from the spectral line shown below.

  \ioaafig{2020_DA_1-f3.pdf}
\item[(e)] (4 points) Calculate the mass of the central BH
  ($M_{\mathrm{vir,BH}}$) in a unit of $M_\odot$.
\end{description}

\probhead{12.16}{Distance to the Large Magellanic Cloud}{2023}{Katowice, Poland}{DA Q1}{50}
% source id: 2023_DA_Q1   tag: eclipsing-binary distance to LMC
In 2019 an international collaboration led by Polish astronomers measured,
with very high precision and accuracy, the distance to the Large Magellanic
Cloud (LMC), a satellite galaxy of the Milky Way. In this way they set the
zero point of the extragalactic distance scale, which allowed for a very
precise measurement of the Hubble constant. Their method involved measuring
the distances to 20 eclipsing binary stars in the LMC, using the concept of
the surface brightness $S_V$ of a star defined as:
\[ S_V = m_V + 5 \log_{10} \theta, \]
where $m_V$ is the magnitude of a star in the optical $V$ band and $\theta$ is
the angular diameter of the star on the sky in milliarcseconds (mas).

The quantity $S_V$ can be understood as the magnitude of a star with an
angular diameter of 1\,mas. An empirical relation has been established between
$S_V$ and the colour index $(m_V - m_K)$, where $m_V$ and $m_K$ are magnitudes
in the $V$-band and infrared $K$-band. This is shown in the figure below for
giant stars of spectral types G and K.

\ioaafig{2023_DA_Q1-f1.pdf}

Using this relation, the distance to an eclipsing binary system can be
determined by deriving the physical radii of the components (using photometry
and spectroscopy), and comparing these with the angular diameters predicted by
the $S_V - (m_V - m_K)$ relation.

The table below gives the parameters of three eclipsing binary stars. $R_1$
and $R_2$ are the radii of each component, $V_{1+2}$ and $K_{1+2}$ are the
total brightness in magnitudes of the binary in the $V$- and $K$-bands, and
$L_2/L_1$ is the luminosity ratio of the components in each band.

\begin{center}
\begin{tabular}{l|c|c|c|c|c|c}
source ID & $R_1$ [$R_\odot$] & $R_2$ [$R_\odot$] & $V_{1+2}$ [mag]
  & $K_{1+2}$ [mag] & $L_2/L_1$ ($V$) & $L_2/L_1$ ($K$) \\
\hline
OGLE-LMC-ECL-03160 & 17.03 & 37.42 & 16.73 & 14.10 & 2.80  & 4.23  \\
OGLE-LMC-ECL-10567 & 24.60 & 36.64 & 16.15 & 13.83 & 1.41  & 1.99  \\
OGLE-LMC-ECL-18365 & 37.30 & 15.94 & 16.27 & 14.01 & 0.206 & 0.188 \\
\end{tabular}
\end{center}

Apply the method outlined above to the three eclipsing binary systems and
calculate the distance to the LMC in kiloparsecs. Estimate the total error of
the result. Assume that the fitting of the $S_V - (m_V - m_K)$ relation
contributes to a bias of up to 0.8\% in all measurements simultaneously.

Hint: in your calculations keep at least three significant figures and two
decimal places. Assume that interstellar extinction is negligible and that the
angular size of the LMC is small.

\clearpage
\section*{Solutions to Chapter 12}

\solhead{12.1}{}
% source id: 2008_TH_S15
Absolute B magnitude of NGC 2639:
\[
M_B = -9.95 \log v_{max} + 3.15 = -9.95 \log(324) + 3.15 = -21.83
\hfill(20\%)
\]
\[
\log R_{25} = -0.249 M_B - 4.00 = 1.4357
\;\;\Longrightarrow\;\;
R_{25} = 27.2678~\mathrm{pc} = 8.41 \times 10^{19}~\mathrm{cm}
\hfill(20\%)
\]
% sic: source says "pc" and "8.41e19 cm"; from the definition R25 is in kpc
% (27.2678 kpc = 8.41e22 cm), and the final mass below is consistent with kpc.
Mass estimation of NGC 2639:
\[
\mathcal{M} = \frac{v^2 R_{25}}{G}
= \frac{(324 \times 10^{5})^{2}\,(8.41 \times 10^{19})}{6.6726 \times 10^{-8}}
= 1.32 \times 10^{42}~\mathrm{kg}
= 6.62 \times 10^{11}~\mathcal{M}_\odot
\hfill(20\%)
\]
% sic: mixed CGS/SI units as in the original
The B magnitude of the Sun:
\[
M_{B\odot} = M_{V\odot} + (m_{B\odot} - m_{V\odot}) = 4.82 + 0.64 = 5.46
\hfill(20\%)
\]
meaning the Sun is fainter in the blue band relative to the A0 standard star.

Luminosity of NGC 2639 in solar luminosity ($L_\odot$):
\[
L = 10^{0.4(M_{B\odot} - M_B)} L_\odot
= 10^{0.4(5.46 + 21.83)} L_\odot
= 8.24 \times 10^{10}\,L_\odot
\hfill(20\%)
\]

\solhead{12.2}{}
% source id: 2011_TH_L2
\begin{description}
\item[1)] The luminosity distance is defined to satisfy a canonical inverse
  square law of the bolometric luminosity $L$ (total emitted radiation power)
  and observed flux $F$ (flux--luminosity relationship)
  \[ F = \frac{L}{4\pi d_L^2} \tag*{(1)} \]
  Units of the flux and luminosity are expressed in SI
  ($\mathrm{W\,m^{-2}}$ and W respectively), or CGS
  ($\mathrm{erg\,s^{-1}\,cm^{-2}}$, $\mathrm{erg\,s^{-1}}$ respectively).
  \hfill(3 points)

\item[2)] Only the part of the spectrum (namely B-band) is observed, but the
  inverse square law is applicable also to the monochromatic luminosity
  $L_\lambda(\lambda)$ and monochromatic flux
  $S(\lambda)$,\footnote{Monochromatic flux is sometimes called as spectral
  density of flux (energy per unit time per unit area per unit wavelength)
  while monochromatic luminosity as spectral density of the luminosity (energy
  per unit time per unit wavelength). Their form is described by SED and thus
  spectral index.} taking into account that observed wavelengths are not equal
  to emitted ones. (Observers frame is not the same as a frame of galaxy.) For
  a small wavelength interval
  \[ S(\lambda_o)\cdot\Delta\lambda_o
     = \frac{L_\lambda(\lambda_e)}{4\pi d_L^2}\cdot\Delta\lambda_e , \]
  Observed fraction of the total spectrum is redshifted into the frame of the
  observer, and a relationship between the observed and emitted wavelengths,
  $\lambda_o$ and $\lambda_e$, can be expressed as
  \[ \lambda_e = \lambda_o/(1+z) . \]
  So,
  \[ S(\lambda)\Delta\lambda = \frac{1}{4\pi d_L^2}\cdot
     L_\lambda\!\left(\frac{\lambda}{1+z}\right)\cdot\frac{1}{1+z}\,
     \Delta\lambda . \tag*{(2)} \]
  \hfill(9 points)

\item[3)] The absolute luminosity is defined as the flux would be observed at
  a distance of 10\,pc from the compact source at rest (not redshifted). It
  means that the frame would be the same as for the object.
  \[ S_{10}(\lambda) = \frac{L_\lambda(\lambda)}{4\pi d_{10}^2}, \tag*{(3)} \]
  where $d_{10}$ is the distance equal to 10\,pc, $S_{10}(\lambda)$ denotes a
  monochromatic flux in the wavelength of $\lambda$ measured at 10\,pc from
  the object.
  \hfill(3 points)

\item[4)] Dividing eq.~(2) by (3) we get
  \[ \frac{S(\lambda)}{S_{10}(\lambda)}
     = \left(\frac{d_{10}}{d_L}\right)^2 \cdot
       \frac{L_\lambda\!\left(\dfrac{\lambda}{1+z}\right)}{L_\lambda(\lambda)}
       \cdot \frac{1}{1+z} \tag*{(4)} \]
  for all the wavelengths.

  The lower edge of validity of SED approximation is equal to 250\,nm. This
  corresponds to $1.5\cdot 250\ \mathrm{nm} = 375$\,nm what is below the blue
  cut-off of the B filter,\footnote{Blue band: Effective Wavelength Midpoint
  $\lambda\mathrm{eff}=445$\,nm, FWHM${}=94$\,nm
  [398\,nm $\div$ 492\,nm]} so SED approximation is applicable. Using the SED
  of the galaxy, we get
  \begin{align*}
    \frac{L_\lambda\!\left(\dfrac{\lambda}{1+z}\right)}{L_\lambda(\lambda)}
      \cdot \frac{1}{1+z}
      &= \left(\frac{1}{1+z}\right)^{4}\frac{1}{1+z}
       = (1+z)^{-5} \tag*{(5)} \\
    \frac{S(\lambda)}{S_{10}(\lambda)}
      &= \left(\frac{d_{10}}{d_L}\right)^{2}\cdot(1+z)^{-5} \tag*{(6)}
  \end{align*}
  The fraction of the monochromatic fluxes $S(\lambda)/S_{10}(\lambda)$ does
  not depend on the wavelength! So the relation (6) is valid also for blue
  band fluxes.
  \hfill(9 points)

\item[5)] The blue fluxes $F$ and $F_{10}$ have to be replaced by the
  corresponding magnitudes. Using Pogson's Ratio it can be obtained:
  \begin{align*}
    m_B - M_B &= -2.5\log\left(\frac{F}{F_{10}}\right)
      = 5\log\left(\frac{d_L}{d_{10}}\right) + 12.5\log(1+z) \\
    M_B &= m_B - 5\log\left(\frac{d_L}{d_{10}}\right)
      + 12.5\log(1+z) \tag*{(7)} % sic: source has "+"; the numerical
      % substitution below (and the first line) require "-"
  \end{align*}
  Substituting numerical data one can obtain:
  \[ M_B = 20.40^{\mathrm{mag}} - 5\log\!\left(2754\cdot 10^{5}\right)
     - 12.5\log(1.5) = -24.00^{\mathrm{mag}} \]
  \[ M_B = -24.00^{\mathrm{mag}} \tag*{(8)} \]
  Please note, that precision is determined by $m_B$!
  \hfill(3 points)

\item[6)] Only the cD galaxies are so luminous, while the normal elliptical
  galaxies are substantially weaker. Accordingly, the galaxy is not a member
  of the cluster, but is a foreground object.
  \hfill(3 points)
\end{description}

\solhead{12.3}{}
% source id: 2012_TH_L2
% Official solution labels these parts B1--B6 (question 2 = section B of the
% 2012 long-question solutions); they correspond to parts 2.1--2.6 of the
% statement. The A section in the same crop belongs to long question 1 and is
% not transcribed here.
\begin{description}
\item[B1.] \mbox{}
  \begin{align*}
    &1^{\mathrm{st}}\ \text{line:}\quad \{(0.0);\,(15.210)\}\ \
      (10^3\,\mathrm{ly},\ \mathrm{km/s}) \\
    &2^{\mathrm{nd}}\ \text{line:}\quad \{(15.210);\,(60.200)\}\ \
      (10^3\,\mathrm{ly},\ \mathrm{km/s})
  \end{align*}
  \[ V(D) = \begin{cases} 14D, & 0 \le D < 15 \\
       213 - 0.22D, & 15 \le D \le 60 \end{cases}, \qquad
     D \text{ in } 10^3\,\mathrm{l.y.} \text{ and } V \text{ in km/s} \]

\item[B2.] Using the equations from (a):
  \[ \omega(D) = \frac{V(D)}{D} \;\Rightarrow\;
     \omega(D) = \begin{cases} 14, & 0 \le D < 15 \\[2pt]
       \dfrac{213}{D} - 0.22, & 15 \le D \le 60 \end{cases}
     \quad \frac{\mathrm{km/s}}{10^3\,\mathrm{l.y.}} \]
  \[ \omega_{\mathrm{wave}}(D) = 2\omega(D) \]
  The time that a spiral arm takes to have one more turn around the galactic
  center:
  \[ \omega_{\mathrm{rel}} = \omega_{\mathrm{wave}}^{\mathrm{max}}
     - \omega_{\mathrm{wave}}^{\mathrm{min}} \]
  \[ \omega_{\mathrm{rel}} = \frac{1}{2}\left[
       \left(\frac{213}{15} - 0.22\right)
       - \left(\frac{213}{80} - 0.22\right)\right]
     = 23.9\,\frac{\mathrm{km/s}}{10^3\,\mathrm{l.y.}}
     % sic: as printed; the prefactor 1/2 and the denominator 80 do not
     % reproduce 23.9 (1/2 x 11.54 = 5.77), which is closer to 2 x 11.54
  \]
  \[ P_{\mathrm{sin}} = \frac{2\pi}{\omega_{\mathrm{rel}}}
     = 8.15\cdot 10^{7}\ \text{years} \]

\item[B3.] Tully--Fisher relation estimates the luminosity of the galaxy, and
  considering $\Delta V = 2V_{\mathrm{max}}$:
  \[ L = 16\,k\,V_{\mathrm{max}}^{4} \]
  We can compare the magnitude of the galaxy with the Sun:
  \[ m - M = 5\log d - 5 \;\Rightarrow\;
     m - \left(M_{\mathrm{Sun}}
       - 2.5\log\frac{L}{L_{\mathrm{Sun}}}\right) = 5\log d - 5 \]
  \[ \log d = 1 + \frac{m - M_{\mathrm{Sun}}
       + 2.5\log\!\left(\dfrac{16\,k\,V_{\mathrm{max}}^{4}}
       {L_{\mathrm{Sun}}}\right)}{5} \]
  From the graph: $V_{\mathrm{max}} = 230$\,km/s:
  \[ \log d = 1 + \frac{8.5 - 4.8
       + 2.5\log\!\left(16\cdot 0.317\cdot 230^{4}\right)}{5}
     = 6.816 \;\Rightarrow\; d = 6.5\,\mathrm{Mpc} \]

\item[B4.] We can estimate the recession velocity of the galaxy using
  Hubble's law:
  \[ V_0 = H\cdot d \]
  The maximum and minimum wavelengths occur where the radial velocity is
  $V_{\mathrm{max}}$, the radial velocity is:
  \[ V = V_0 \pm V_{\mathrm{max}} \]
  Using $\dfrac{\Delta\lambda}{\lambda} = \dfrac{v}{c}$, the wavelengths are:
  \begin{align*}
    \lambda_{\mathrm{max}} &= \lambda\left(1
      + \frac{Hd + V_{\mathrm{max}}}{c}\right) \\
    \lambda_{\mathrm{min}} &= \lambda\left(1
      + \frac{Hd - V_{\mathrm{max}}}{c}\right)
  \end{align*}
  Calculating the values, using $\lambda = 656.28$\,nm:
  \begin{align*}
    \lambda_{\mathrm{max}} &= 657.79\,\mathrm{nm} \\
    \lambda_{\mathrm{min}} &= 656.78\,\mathrm{nm}
  \end{align*}

\item[B5.] Supposing a mass with spherical symmetry distribution and a
  circular orbit, we have:
  \[ \frac{V^2}{D} = \frac{G M_{\mathrm{int}}}{D^2} \;\Rightarrow\;
     M_{\mathrm{int}} = \frac{V^2 D}{G} \]
  From the graph $V(3.0\cdot 10^{4}\ \mathrm{l.y.}) = 225$\,km/s, So the mass
  is:
  \[ M_{\mathrm{int}} = 2.15\cdot 10^{41}\,\mathrm{kg}
     = 1.08\cdot 10^{11}\,M_{\mathrm{sun}} \]

\item[B6.] As the mass is derived from the rotation curve, it is necessary to
  take into account that part of the mass is dark matter, and only part of the
  mass is in the stars ($\rho_{\mathrm{star}} = 1/3$). The fraction of
  baryonic matter is:
  \[ \rho_{\mathrm{matter}} = \frac{f_{\mathrm{baryon}}}{f_{\mathrm{matter}}}
     = \frac{4}{22+4} = \frac{4}{26} \]
  The number of stars is then:
  \[ N = \rho_{\mathrm{matter}}\cdot\rho_{\mathrm{star}}\,
     \frac{M_{\mathrm{galaxy}}}{M_{\mathrm{Sun}}}
     = \frac{4}{26}\,\frac{1}{3}\,1.08\cdot 10^{11}
     = 5.5\cdot 10^{9}\ \text{stars} \]
\end{description}

\solhead{12.4}{}
% source id: 2015_TH_S7
We have the total energy of the escaping galaxy:
\[
E_T = E_K + E_P = \frac{mV^2}{2} - \frac{GMm}{R}
\hfill(15)
\]
where $V$ is the velocity of the galaxy and $R$ the radius of the galaxy
cluster. The galaxy will escape from the galaxy cluster if the total energy
$E_T = 0$. \hfill(10)
\[
E_T = 0 \;\rightarrow\; \frac{mV^2}{2} - \frac{GMm}{R} = 0
\quad\Longrightarrow\quad
V^2 = \frac{2GM}{R}
\hfill(15)
\]
Escape velocity:
\[
V^2 = \frac{2G\left(\dfrac{4\pi}{3}R^3\rho\right)}{R}
\hfill(20)
\]
Density of cluster:
\[
\rho = \frac{3V^2}{8\pi G R^2}
\hfill(15)
\]
With $V = 700$~km/s, $R = 10$~Mpc,
$G = 6.6378 \times 10^{-11}~\mathrm{N\,m^2\,kg}$:
% sic: G value and units as in the original (should be 6.674e-11 N m^2 kg^-2)
\[
\rho = 9.20121 \times 10^{-27}~\mathrm{kg/m^3}
\hfill(25)
\]

\solhead{12.5}{Mass of the Local Group}
% source id: 2017_TH_T11
\begin{description}
\item[a)]
  \begin{equation*}
  \frac{1}{2}m_1 v_1^2 + \frac{1}{2}m_2 v_2^2 - \frac{Gm_1 m_2}{r_1 + r_2} = E.
  \tag*{[5]}
  \end{equation*}
  Note that the total energy is a negative quantity.

\item[b)] From $m_1 r_1 = m_2 r_2$ and $r_1 + r_2 = r$, we have
  \[
  r_1 = \frac{m_2}{m_1 + m_2}\,r
  \qquad\text{and}\qquad
  r_2 = \frac{m_1}{m_1 + m_2}\,r.
  \]
  Substituting them into the solution obtained in a), we get
  \begin{align*}
  \frac{1}{2}\mu v^2 - \frac{G(m_1 + m_2)\mu}{r} &= E \\
  v^2 - \frac{2GM}{r} &= \frac{2E}{\mu}
  \end{align*}
  Therefore,
  \begin{equation*}
  v^2 = 2\left(\frac{GM}{r} + \frac{E}{\mu}\right)
  \tag*{[10]}
  \end{equation*}

\item[c)] From b), we have
  \begin{align*}
  v^2 &= \frac{2GM}{r} + \frac{2E}{\mu}
  = (2GM)\left(\frac{1}{r} + \frac{E}{GM\mu}\right) \\
  &= (2GM)\left(\frac{1}{r} - \frac{1}{r_0}\right),
  \quad\text{where } r_0 = -\frac{GM\mu}{E}.
  \tag*{[5]}
  \end{align*}

\item[d)] Substituting $r(\theta)$ in the equation of part c):
  \[
  v^2 = 2GM\left(\frac{1}{r} - \frac{1}{r_0}\right)
  = 2GM\left(\frac{1}{\tfrac{1}{2}r_0(1-\cos\theta)} - \frac{1}{r_0}\right)
  = \frac{2GM}{r_0}\left(\frac{1+\cos\theta}{1-\cos\theta}\right)
  \]
  \[
  v^2 = \frac{2GM}{r_0}\,\frac{(1-\cos^2\theta)}{(1-\cos\theta)^2}
  = \frac{2GM}{r_0}\,\frac{\sin^2\theta}{(1-\cos\theta)^2}
  \;\Longrightarrow\;
  v(\theta) = \left(\frac{2GM}{r_0}\right)^{1/2}\frac{\sin\theta}{1-\cos\theta}
  \]
  Combining $v(\theta)$, $t(\theta)$, and $r(\theta)$:
  \[
  \frac{v(\theta)\,t(\theta)}{r(\theta)}
  = \frac{\left(\dfrac{2GM}{r_0}\right)^{1/2}
      \dfrac{\sin\theta}{1-\cos\theta}\cdot
      \left(\dfrac{r_0^3}{8GM}\right)^{1/2}(\theta - \sin\theta)}
      {\dfrac{r_0}{2}\,(1-\cos\theta)}
  = \frac{2}{r_0}\left(\frac{2GM r_0^3}{r_0\, 8GM}\right)^{1/2}
    \frac{(\sin\theta)(\theta - \sin\theta)}{(1-\cos\theta)^2}
  \]
  \[
  \frac{vt}{r} = \frac{(\sin\theta)(\theta - \sin\theta)}{(1-\cos\theta)^2}
  \]

  \emph{Alternative solution:}
  From $r(\theta) = \dfrac{r_0}{2}(1-\cos\theta)$, we have
  \[
  v = \left(\frac{r_0}{2}\sin\theta\right)\dot\theta
  \]
  From $t(\theta) = \left(\dfrac{r_0^3}{8GM}\right)^{1/2}(\theta-\sin\theta)$,
  we have
  \[
  1 = \left(\frac{r_0^3}{8GM}\right)^{1/2}(1-\cos\theta)\,\dot\theta
  \]
  \begin{equation*}
  \therefore\;
  \frac{vt}{r}
  = \frac{\left(\dfrac{r_0}{2}\sin\theta\right)
      \left(\dfrac{r_0^3}{8GM}\right)^{1/2}(\theta-\sin\theta)}
      {\left(\dfrac{r_0^3}{8GM}\right)^{1/2}(1-\cos\theta)\,
       \dfrac{r_0}{2}(1-\cos\theta)}
  = \frac{(\sin\theta)(\theta-\sin\theta)}{(1-\cos\theta)^2}.
  \tag*{[10]}
  \end{equation*}

\item[e)]
  \begin{align*}
  v &= -118~\mathrm{km\,s^{-1}} = -118 \times 10^{3}~\mathrm{m\,s^{-1}} \\
  r &= 710~\mathrm{kpc} = 2.19 \times 10^{22}~\mathrm{m} \\
  t &= 13700 \times 10^{6}~\mathrm{years} = 4.3233 \times 10^{17}~\mathrm{s}
  \end{align*}
  \begin{equation*}
  \therefore\;
  \frac{vt}{r}
  = -\frac{118 \times 10^{3} \times 4.3233 \times 10^{17}}
          {2.1908 \times 10^{22}}
  = -2.33
  = \frac{(\sin\theta)(\theta-\sin\theta)}{(1-\cos\theta)^2}
  \tag*{[5]}
  \end{equation*}
  The negative value of the left-hand side of this equation implies that
  $\theta$ is greater than $\pi$.

  \begin{center}
  \begin{tabular}{cc}
  $\theta$ & $(\sin\theta)(\theta-\sin\theta)/(1-\cos\theta)^2$ \\
  \hline
  $200^\circ = 3.49$ radians & $-0.348$ \\
  $210^\circ = 3.67$ radians & $-0.598$ \\
  $240^\circ = 4.19$ radians & $-1.95$ \\
  $244^\circ = 4.26$ radians & $-2.24$ \\
  $245^\circ = 4.28$ radians & $-2.32$ \\
  $246^\circ = 4.29$ radians & $-2.40$ \\
  $250^\circ = 4.36$ radians & $-2.77$ \\
  \end{tabular}
  \end{center}
  This gives the value of $\theta = 245^\circ = 4.28$ radians. \hfill[5]

\item[f)] From $r(\theta) = \dfrac{r_0}{2}(1-\cos\theta)$, we have
  $r_{\mathrm{max}} = r_0$, and hence
  \begin{equation*}
  r_{\mathrm{max}} = \frac{2r(\theta)}{1-\cos\theta}
  = \frac{2 \times 710~\mathrm{kpc}}{1 - \cos 245^\circ}
  = 998~\mathrm{kpc}
  \tag*{[5]}
  \end{equation*}
  (This is the maximum separation between the Milky-Way Galaxy and Andromeda
  Galaxy; they started to move towards each other afterwards.)

  The value of $M$ can be calculated from the relation
  $t(\theta) = \left(\dfrac{r_0^3}{8GM}\right)^{1/2}(\theta - \sin\theta)$:
  \[
  M = \left(\frac{r_{\mathrm{max}}^3}{8G}\right)
      \left\{\frac{\theta - \sin\theta}{t(\theta)}\right\}^{2}.
  \]
  From the maximum separation of
  $r_{\mathrm{max}} = 998~\mathrm{kpc} = 3.0795 \times 10^{22}$~m,
  $\theta = 245^\circ = 4.276$~radians,
  $t = t(\theta) = 13700 \times 10^{6}~\mathrm{years}
     = 4.3233 \times 10^{17}$~s, and
  $G = 6.67 \times 10^{-11}~\mathrm{N\,m^2\,kg^{-2}}$, we obtain
  \begin{equation*}
  \therefore\; M = 7.86 \times 10^{42}~\mathrm{kg}
  = \frac{7.86 \times 10^{42}}{1.99 \times 10^{30}}\,M_\odot
  = 3.95 \times 10^{12}\,M_\odot
  \tag*{[5]}
  \end{equation*}
  The estimated number of stars, most of which are of less than a solar mass,
  for Andromeda and Our Galaxy are $4 \times 10^{11}$ and $2 \times 10^{11}$
  respectively. This implies that most of mass of the system is DARK.
\end{description}

\medskip
\noindent\textbf{Note that}
\[
v = \frac{d}{dt}r = +(2GM)^{1/2}\left(\frac{r_0 - r}{r_0 r}\right)^{1/2}
\quad\text{for } r_0 > r
\]
\[
\int (2GM)^{1/2}\,dt = +\int \left(\frac{r_0 r}{r_0 - r}\right)^{1/2} dr
\]
Putting $\dfrac{r}{r_0} \equiv \sin^2\phi$,
\[
t = \left(\frac{r_0^3}{8GM}\right)^{1/2}\{2\phi - \sin 2\phi\} + D
\]
We choose the initial condition such that $m_1$ and $m_2$ are close together,
that is that $r = 0$ at $t = 0$. This implies that $\phi = 0$ at $t = 0$,
hence $D$ must be zero.
\[
t = \left(\frac{r_0^3}{8GM}\right)^{1/2}(\theta - \sin\theta),
\quad \theta \equiv 2\phi.
\]
And from $r = r_0 \sin^2\phi = \dfrac{1}{2}r_0(1 - \cos 2\phi)$, we have:
\[
r = \frac{1}{2}r_0(1 - \cos\theta).
\]
These form the parametric solution:
$r_{\mathrm{max}} = r(\theta = \pi) = r_0$.

\solhead{12.6}{Galactic Outflow}
% source id: 2017_TH_T9
\begin{description}
\item[a)] The systemic or recessional velocity of IRAS 0833+6517 can be
  determined by the Hubble's law
  \begin{align*}
  v_{\mathrm{sys}} = H_0 d
  &= 67.8~\mathrm{km\,s^{-1}\,Mpc^{-1}} \times 80.2~\mathrm{Mpc} \\
  &= 5437.6~\mathrm{km\,s^{-1}}
  \tag*{[2]}
  \end{align*}
  The observed radial velocity can be written in terms of systemic velocity
  and rotational velocity of a galaxy as
  \begin{equation*}
  v_{\mathrm{rad}} = v_{\mathrm{sys}} + v_{\mathrm{rot}} \sin i,
  \tag*{[1]}
  \end{equation*}
  where $v_{\mathrm{rad}}$, $v_{\mathrm{sys}}$, and $v_{\mathrm{rot}}$ are the
  radial, systemic, and rotational velocity of the galaxy and $i$ is an
  inclination angle.

  The rotational velocity of the galaxy can thus be obtained by
  \begin{equation*}
  v_{\mathrm{rot}} = \frac{v_{\mathrm{rad}} - v_{\mathrm{sys}}}{\sin i}
  = \frac{(5850 - 5437.6)~\mathrm{km\,s^{-1}}}{0.39}
  = 1055~\mathrm{km\,s^{-1}}
  \tag*{[2]}
  \end{equation*}
  Answers between 1050 to 1060 are acceptable for full marks.

\item[b)] \textbf{Method 1:} In order to determine the escape velocity, we
  need to know the dynamical mass of the system. In case of the rotation
  dominated galaxy, a dynamical mass can be estimated as
  \begin{equation*}
  M_{\mathrm{dyn}}(<R) = v_{\mathrm{rot}}^2 R / G,
  \tag*{[2]}
  \end{equation*}
  where $M_{\mathrm{dyn}}(<R)$ is the dynamical mass within the radius $R$.
  So the dynamical mass of IRAS 0833+6517 within the radius of 7.8~kpc is
  \begin{align*}
  M_{\mathrm{dyn}}(<R)
  &= \frac{(1055 \times 10^{3}~\mathrm{m\,s^{-1}})^{2}
     \times 7.8 \times 10^{3}~\mathrm{pc} \times 206265~\mathrm{au/pc}
     \times 1.50 \times 10^{11}~\mathrm{m/au}}
     {6.67 \times 10^{-11}~\mathrm{m^{3}\,kg^{-1}\,s^{-2}}} \\
  &= 4.0 \times 10^{42}~\mathrm{kg}
  \tag*{[2]}
  \end{align*}
  or $M_{\mathrm{dyn}}(<R) = 2.0 \times 10^{12}\,M_\odot$.

  Assuming the virial theorem, the escape velocity can be written as
  \begin{equation*}
  v_{\mathrm{esc}} = \sqrt{\frac{2GM}{R}}
  \tag*{[2]}
  \end{equation*}
  Therefore, the escape velocity is
  \begin{align*}
  v_{\mathrm{esc}}
  &= \sqrt{\frac{2 \times 6.67 \times 10^{-11}~\mathrm{m^{3}\,kg^{-1}\,s^{-2}}
      \times 4.0 \times 10^{42}~\mathrm{kg}}
     {7.8 \times 10^{3}~\mathrm{pc} \times 206265~\mathrm{au/pc}
      \times 1.50 \times 10^{11}~\mathrm{m/au}}} \\
  &= 1492~\mathrm{km\,s^{-1}}
  \tag*{[3]}
  \end{align*}

  \textbf{Method 2:} In order to determine the escape velocity, we need to
  know the dynamical mass of the system. In case of the rotation dominated
  galaxy, a dynamical mass can be estimated as
  \begin{equation*}
  M_{\mathrm{dyn}}(<R) = v_{\mathrm{rot}}^2 R / G,
  \tag*{[2]}
  \end{equation*}
  where $M_{\mathrm{dyn}}(<R)$ is the dynamical mass within the radius $R$.
  According to the virial theorem, the escape velocity is
  \begin{equation*}
  v_{\mathrm{esc}} = \sqrt{\frac{2GM}{R}}
  \tag*{[2]}
  \end{equation*}
  Substituting the dynamical mass into the above equation, we get
  \begin{align*}
  v_{\mathrm{esc}} &= \sqrt{2\,\frac{GM}{R}} = \sqrt{2}\,v_{\mathrm{rot}}
  \tag*{[3]} \\
  &= \sqrt{2}\,(1055~\mathrm{km\,s^{-1}}) = 1492~\mathrm{km\,s^{-1}}
  \tag*{[2]}
  \end{align*}
  Answers between 1487 to 1497 are acceptable for full marks.

\item[c)] To examine whether or not outflowing gas can escape from the
  galaxy, one needs to derive the outflow velocity of the gas from the
  blueshifted absorption line.

  The outflow velocity or the velocity offset of any emission or absorption
  lines can be calculated by the following equation:
  \begin{equation*}
  \Delta v = \frac{\Delta\lambda}{\lambda}\,c
  \tag*{[2]}
  \end{equation*}
  \ioaafig{2017_TH_T9-s1.pdf}
  For the C~II $\lambda 1335$ absorption line (Figure~1), we obtain
  \begin{align*}
  \Delta v_{\mathrm{obs}}(\mathrm{C~II}\,\lambda 1335)
  &= 299792.458~\mathrm{km\,s^{-1}}
     \left(\frac{1332 - 1335}{1335}\right) \\
  &= -673.7~\mathrm{km\,s^{-1}}
  \tag*{[2]}
  \end{align*}
  Answers between $-300~\mathrm{km\,s^{-1}}$ to $-900~\mathrm{km\,s^{-1}}$
  are acceptable with full marks.

  The velocity offsets of the C~II $\lambda 1335$ absorption line is
  $-673.7~\mathrm{km\,s^{-1}}$. It means that the outflowing gas is going out
  of the galaxy at this speed, which is smaller than the escape velocity
  obtained in (b).

  We conclude that the outflowing gas CANNOT escape from
  IRAS 0833+6517. \hfill[2]
\end{description}

\solhead{12.7}{A Possible Dark Matter Deficient Galaxy}
% source id: 2018_TH_T8
\begin{description}
\item[(1)] Total area within the ellipse (actually, half-light ellipse) of
  NGC 1052-DF2 is:
  \begin{equation*}
  A_{DF2} = \pi a b = 0.85\pi \times 22.6''^{\,2}
  = 1363.9~\mathrm{arcsec^2}
  \tag*{(2 points)}
  \end{equation*}
  Magnitude of the part of the galaxy within the ellipse:
  \begin{equation*}
  m_{ell} = SB - 2.5 \times \log(A_{DF2})
  = 24.7 - 2.5 \times \log(1363.9) = 16.9
  \tag*{(2 points)}
  \end{equation*}
  Total magnitude of the whole galaxy:
  \begin{equation*}
  m_{gal} = m_{ell} - 2.5 \times \log(2) = 16.1
  \tag*{(2 points)}
  \end{equation*}
\item[(2)] Absolute mag of the galaxy:
  \begin{equation*}
  M_{0,gal} = m_{gal} + 5 - 5 \times \log(D) = -15.4
  \tag*{(2 points)}
  \end{equation*}
  Convert to solar luminosity:
  \begin{equation*}
  \frac{L_{gal}}{L_\odot} = 10^{-0.4(M_{0,gal} - M_{0,\odot})}
  = 1.2 \times 10^{8}
  \tag*{(2 points)}
  \end{equation*}
  Thus, the total stellar mass should be:
  \begin{equation*}
  \frac{M_{gal}}{M_\odot} = \frac{L_{gal}}{L_\odot} \times \frac{M}{L}
  = 2.4 \times 10^{8}
  \tag*{(2 points)}
  \end{equation*}
\item[(3)] Mean galactocentric distance of globular clusters:
  \begin{equation*}
  r_{gc} = \frac{\theta D}{206265} = 7.6~\mathrm{kpc}
  \tag*{(2 points)}
  \end{equation*}
  Using viral theorem: % sic: "viral" for "virial" as in the original
  \begin{gather*}
  \langle 2K \rangle + \langle U \rangle = 0, \\
  \langle K \rangle = \tfrac{1}{2}M\langle v^2 \rangle
  = \tfrac{1}{2}M\sigma^2, \\
  \langle U \rangle = \frac{3}{5}\,\frac{GM^2}{r}
  \tag*{(4 points)}
  \end{gather*}
  Thus the maximal dynamical mass of this system should be:
  \begin{equation*}
  M_{dyn} = \frac{5 r_{gc}\sigma^2}{3G} = 2.1 \times 10^{8}\,M_\odot,
  \tag*{(3 points)}
  \end{equation*}
  even less than the total stellar mass within this radius.
\item[(4)] From Questions 2 and 3, we know:
  \[
  M_{star} \propto L_{gal} \propto F_{gal} \times D^2
  \]
  However,
  \begin{equation*}
  M_{dym} \propto D % sic: "dym" as in the original
  \tag*{(2 points)}
  \end{equation*}
  Thus
  \begin{equation*}
  M_{dym}/M_{star} \propto 1/D
  \tag*{(2 points)}
  \end{equation*}
  If the distance measured is smaller by a certain factor, the stellar mass to
  dynamical mass ratio would increase by the same factor, thus the dark matter
  in NGC 1052-DF2 would not be as deficient as Dragonfly team claimed.
\end{description}

\solhead{12.8}{X-ray emission from galaxy clusters}
% source id: 2023_TH_Q11
\begin{description}
\item[(a)] Concentrations of electrons, $N_e$, and He nuclei, $N_{He}$, is
  related to the concentration of H nuclei, $N_H$:
  \begin{equation*}
  N_{He} = 0.1\,N_H
  \tag*{(1 point)}
  \end{equation*}
  \begin{equation*}
  N_e = N_H + 2N_{He} = 1.2\,N_H
  \tag*{(2 points)}
  \end{equation*}
  The total X-ray emission is a sum of the bremmstrahlung generated by
  % sic: "bremmstrahlung" as in the original
  interaction of electrons with H and He nuclei:
  \[
  L = 6 \cdot 10^{-41}\, N_e\, T^{\frac{1}{2}} V \left(N_H + Z_{He}^2 N_{He}\right),
  \]
  where $Z_{He} = 2$. The luminosity $L$ is expressed by the concentration of
  H nuclei:
  \begin{gather*}
  L = 6 \cdot 10^{-41}\; 1.2\,N_H\; T^{\frac{1}{2}}\, V\; 1.4\,N_H, \\
  \Longrightarrow\;
  L = (10.08 \times 10^{-41})\, T^{\frac{1}{2}} V N_H^2
  \approx 10^{-40}\, T^{\frac{1}{2}} V N_H^2.
  \tag*{(3 points)}
  \end{gather*}
  The volume $V$ is given by:
  \begin{equation*}
  V = \frac{4}{3}\pi R^3 = 1.54 \cdot 10^{67}~\mathrm{m^3}.
  \tag*{(3 points)}
  \end{equation*}
  Thus, the concentration of $N_H$:
  \begin{equation*}
  N_H = \left(\frac{L}{10^{-40}\, T^{\frac{1}{2}} V}\right)^{1/2}
  \approx 8.48 \times 10^{2}~\mathrm{m^{-3}}.
  \tag*{(3 points)}
  \end{equation*}
  To obtain the total mass of the plasma, $M$, one should multiply the volume
  $V$ by the the sum of H and He mass densities:
  % sic: "the the" as in the original
  \begin{equation*}
  M = V\left(N_H m_H + N_{He} m_{He}\right)
  = 3.03 \times 10^{43}~\mathrm{kg}
  \approx 1.52 \times 10^{13}\,M_\odot,
  \tag*{(4 points)}
  \end{equation*}
  where $m_H$ and $m_{He}$ are the masses of hydrogen and helium atoms.

\item[(b)] Let $L$ be the mean free path of a photon. The number of electrons
  in a cylinder with a cross section area of $\sigma$ and a length of $L$
  equals
  \[
  N = nL\sigma = 1 \;\Longrightarrow\; L = 1/(n\sigma).
  \]
  Assuming $n = N_e = 1.2 N_H = 1.018 \times 10^{3}~\mathrm{m^{-3}}$, we get
  \begin{equation*}
  L = \frac{1}{1.018 \times 10^{3} \cdot 6.65 \times 10^{-29}}
  = 1.48 \times 10^{25}~\mathrm{m}
  = 4.79 \times 10^{8}~\mathrm{pc}
  \approx 500~\mathrm{Mpc}.
  \tag*{(5 points)}
  \end{equation*}
  (The radius of the cluster is $\sim 1000$ times smaller than $L$, so the
  interactions between the CMB photons and hot electrons are rare.)

\item[(c)] Using Wien's law,
  \begin{gather*}
  \lambda = \frac{b}{T} = \frac{2.898 \times 10^{-3}}{2.73}
  = 1.07 \times 10^{-3}~\mathrm{m} \\
  \Longrightarrow\; E_0 = \frac{hc}{\lambda} = 1.85 \times 10^{-22}~\mathrm{J}
  \tag*{(4 points)}
  \end{gather*}
  % sic: parts (c)+(d) point values in the source solution (4+3+2) do not
  % match the statement's split (3+6); transcribed as printed.

\item[(d)] We first have to estimate the typical velocity $v_e$ of electrons
  using the formula for the kinetic energy of the particles in a gas.
  \[
  \frac{1}{2}mv^2 = \frac{3}{2}kT
  \;\Longrightarrow\;
  v_e = \sqrt{\frac{3kT}{m_e}}
  \]
  \begin{equation*}
  \therefore\; v_e
  = \sqrt{\frac{3 \cdot 1.381 \times 10^{-23} \cdot 8 \times 10^{7}}
               {9.109 \times 10^{-31}}}
  = 6.032 \times 10^{7}~\mathrm{m\,s^{-1}} = 0.202c
  \tag*{(3 points)}
  \end{equation*}
  The energy of upscattered photons is:
  \begin{equation*}
  E' = \frac{1+\beta}{1-\beta}\,E_0 \approx 1.5 E_0
  = 2.78 \times 10^{-22}~\mathrm{J}.
  \tag*{(2 points)}
  \end{equation*}
  \emph{Note:} The formula for the kinetic energy of particles in an ideal gas
  remains valid even at such high temperatures. This is because
  \[
  x = \frac{m_e c^2}{kT}
  = \frac{9.109 \times 10^{-31} \cdot (2.998 \times 10^{8})^2}
         {1.381 \times 10^{-23} \cdot 8 \times 10^{7}}
  = 74.1 \gg 1,
  \]
  and the electrons can be treated as non-relativistic. The full relativistic
  formula for the rms velocity of particles in an ideal gas is
  \[
  \langle v_e^2 \rangle = \frac{xc^2}{K_2(x)}
  \int_0^\infty \frac{\sinh^4\varphi}{\cosh\varphi}\,
  e^{-x\cosh\varphi}\,d\varphi,
  \]
  where $K_2(x)$ is a modified Bessel function of the second kind. It can be
  shown that for $x = 74.1$ this formula yields
  $\sqrt{\langle v_e^2 \rangle} = 0.198c$, which is virtually identical to the
  prediction of the non-relativistic formula.
\end{description}

\solhead{12.9}{Galaxy Cluster}
% source id: 2024_TH_T2
If $U$ and $K$ are the total gravitational potential energy and the total
kinetic energy of the cluster, respectively, the virial theorem says that:
\begin{equation*}
U + 2K = 0
\tag*{(2.0)}
\end{equation*}
Let's say that there are $N$ galaxies of masses $m_i$, where $i = 1,\dots,N$,
moving with velocities $\vec{u}_i = \vec{v}_i - \vec{v}_0$ with respect to the
cluster, where $\vec{v}_0$ is the velocity of the cluster as a whole.
Therefore, in the cluster frame, $K$ is given by
\[
K = \sum_{i=1}^{N} \frac{1}{2} m_i u_i^2
\]
As an estimation, consider that all galaxies have mass equal to the average
mass $\langle m \rangle = M_T/N$:
\begin{equation*}
K = \frac{1}{2}\,\frac{M_T}{N} \sum_{i=1}^{N} u_i^2
= \frac{1}{2} M_T\,\frac{\sum u_i^2}{N}
= \frac{1}{2} M_T \sigma_v^2
\tag*{(2.0)}
\end{equation*}
where $\sigma_v = \sqrt{\langle u_i^2 \rangle}$ is the root-mean-square galaxy
speed, also known as velocity dispersion. Notice, however, that only the
dispersion in redshift is given, which translates into a radial velocity
dispersion:
\begin{equation*}
\sigma_{v_r} = c \cdot \sigma_z
= 2.998 \times 10^{8}~\mathrm{m/s} \times 0.0005
= 1.499 \times 10^{5}~\mathrm{m/s}
\tag*{(1.0)}
\end{equation*}
For completeness of the solution, we relate $\sigma_v$ and $\sigma_{v_r}$ by
assuming three-dimensional isotropy for velocities, which yields
$\sigma_v^2 = 3\sigma_{v_r}^2$. $K$ is then given by:
\[
K = \frac{3}{2} M_T \sigma_{v_r}^2
\]
The gravitational potential energy, on the other hand, is given by:
\[
U = -\sum_{i \neq j} \frac{G m_i m_j}{r_{ij}}
\]
This can be estimated in different ways. Usually, the spatial distribution of
the large number of galaxies may be approximated as uniform, so that one
estimates $U$ as the binding energy of a homogeneous spherical mass
distribution, which is given by:
\begin{equation*}
U = -\frac{3}{5}\,\frac{G M_T^2}{R}
\tag*{(2.0)}
\end{equation*}
And, assuming the cluster is spherical, its radius is:
\begin{equation*}
R = \frac{D_A \cdot \Delta\theta}{2}
= \frac{1500 \times 10^{6} \times 206265 \times 1.496 \times 10^{11}
        \times 10^{-2}}{2}
= 2.3143 \times 10^{23}~\mathrm{m}
\tag*{(1.0)}
\end{equation*}
Now, using the virial theorem, it is possible to calculate the total mass:
\begin{align*}
&\therefore\;
-\frac{3}{5} G M_T^2\,\frac{1}{R} + 2 \cdot \frac{3}{2} M_T \sigma_{v_r}^2 = 0
\;\Rightarrow\;
M_T = \frac{5 R \sigma_{v_r}^2}{G} \\
&M_T = \frac{5 \times 2.314 \times 10^{23}
             \times (1.499 \times 10^{5})^2}{6.67 \times 10^{-11}}~\mathrm{kg}
\\
&M_T \approx 3.9 \times 10^{44}~\mathrm{kg} = 2.0 \times 10^{14}\,M_\odot
\tag*{(2.0)}
\end{align*}

\medskip
\noindent\emph{Appendix on calculation of $U$:} Other valid methods for
estimating $U$ include explicitly writing the sum as:
\begin{align*}
U &= -\frac{1}{2} G \left(\sum_i m_i\right)^{2}
     \left\langle \frac{1}{r_{ij}} \right\rangle \\
U &= -\frac{1}{2}\,\frac{G M_T^2}{R}
\end{align*}
where the factor of $1/2$ comes from counting all the unique pairs, and
$\langle r_{ij}^{-1} \rangle$ was taken (as an approximation) as $R^{-1}$.
Or, alternatively, using dimensional analysis to argue that:
\[
U = -\frac{G M_T^2}{R}.
\]

\solhead{12.10}{Quasars}
% source id: 2025_TH_T05
\begin{description}
\item[(T05.1)] Angular diameter distance $= D_L/(1+z)^2$
  \begin{align*}
  \text{Scale} &= \text{Angular diameter distance} \times
    \frac{\pi}{60 \times 60 \times 180}~\mathrm{ly/arcsecond} \\
  &= 20.7~\mathrm{ly~milliarcsecond^{-1}}.
  \tag*{(1.0)}
  \end{align*}
  \begin{center}
  \begin{tabular}{cccc}
  Date of observation & $\phi_{\mathrm{blob}}$ (milliarcsecond)
    & $l_{\mathrm{blob}}$ (ly) & $\beta_{\mathrm{app}}$ \\
  \hline
  $T_0$            & 1.6 & 33 & $\times$ \\
  $T_0 + 605$~d    & 2.1 & 44 & 6.6 \\
  $T_0 + 1186$~d   & 2.3 & 48 & 2.5 \\
  $T_0 + 1786$~d   & 2.8 & 58 & 6.1 \\
  $T_0 + 2338$~d   & 3.1 & 64 & 4.0 \\
  \end{tabular}
  \hfill(3.0)
  \end{center}
  The average apparent velocity can be calculated by using the distance and
  time between the first and the last point which yields
  $\beta^{\mathrm{ave}}_{\mathrm{app}} = 4.8$.

\item[(T05.2a)] Radiation from the core always takes the same time to reach
  us, but the time taken by radiation from the moving blob in the jet
  decreases since its distance from us reduces.
  \ioaafig{2025_TH_T05-s1.pdf}

  As shown in the sketch above, a signal emitted from the moving blob at B
  covers a distance $(\beta c)\Delta t \cos\theta$ less than a signal emitted
  at A.

  Thus the apparent time difference between the light received by the
  observer is
  \begin{align*}
  \Delta t_{\mathrm{app}} = t_B - t_A
  &= \left(t_0 + \Delta t - \frac{(\beta c)\Delta t \cos\theta}{c}\right)
     - t_0 \\
  &= \Delta t\,(1 - \beta\cos\theta) \\
  \Longrightarrow\;
  \frac{\Delta t_{\mathrm{app}}}{\Delta t}
  &= (1 - \beta\cos\theta).
  \tag*{(2.0)}
  \end{align*}

\item[(T05.2b)] The apparent velocity on the sky is
  \begin{align*}
  \beta_{\mathrm{app}} = \frac{v_{\mathrm{app}}}{c}
  &= \frac{\beta\,\Delta t \sin\theta}{\Delta t_{\mathrm{app}}} \\
  \beta_{\mathrm{app}}
  &= \frac{\beta\sin\theta}{1 - \beta\cos\theta}
  \tag*{(2.0)}
  \end{align*}
  The resulting formula indicates that the measured apparent velocity of the
  jet's component is increased for relativistic features and can be higher
  than speed of light for small values of angle $\theta$.

\item[(T05.3a)] For $\beta_{\mathrm{app}} = 1$, we obtain
  \[
  \beta(\sin\theta + \cos\theta) = 1
  \;\Longrightarrow\;
  \beta = \frac{1}{\sin\theta + \cos\theta}
  \]
  Superluminal motion occurs if $\beta_{\mathrm{app}} > 1$; then
  \[
  \beta(\sin\theta + \cos\theta) > 1.
  \]
  This defines the superluminal region which is to be shaded.
  \ioaafig{2025_TH_T05-s2.pdf}

\item[(T05.3b)] It may be seen from graph drawn in (T05.3a) that the minimum
  value of $\beta$ for which superluminal motion is apparent occurs at
  $\theta = 45^\circ$, i.e., $\beta = 1/\sqrt{2}$.
  \begin{equation*}
  \beta_{\mathrm{low}} = 1/\sqrt{2} \approx 0.707,
  \tag*{(1.0)}
  \end{equation*}
  and the corresponding orientation angle,
  $\theta_{\mathrm{low}} = 45^\circ$. \hfill(1.0)

\item[(T05.4)] We know
  \[
  \beta_{\mathrm{app}} = \frac{\beta\sin\theta}{1 - \beta\cos\theta}
  \]
  Thus
  \[
  \beta\left(\frac{\sin\theta + \beta_{\mathrm{app}}\cos\theta}
                  {\beta_{\mathrm{app}}}\right) = 1
  \]
  Since $\beta \le 1$, we have
  \begin{equation*}
  \frac{\sin\theta + \beta_{\mathrm{app}}\cos\theta}{\beta_{\mathrm{app}}} > 1
  \tag*{(1.0)}
  \end{equation*}
  or
  \[
  \beta_{\mathrm{app}} \tan(\theta/2) < 1.
  \]
  Therefore, the maximum viewing angle for a jet with a given
  $\beta_{\mathrm{app}}$ is
  \begin{equation*}
  \theta_{\mathrm{max}}
  = 2\tan^{-1}\!\left(\frac{1}{\beta_{\mathrm{app}}}\right)
  \tag*{(1.0)}
  \end{equation*}

\item[(T05.5)] Time-scale $\delta t$ of any substantial variability from a
  source cannot be shorter than the light crossing time of the source.

  This then gives an upper limit to the size $\lesssim c\,\delta t$, where $c$
  is the speed of light. Therefore,
  \begin{equation*}
  c\,\delta t \gtrsim 5 R_{SC} = 5 \times \frac{2GM}{c^2}.
  \tag*{(1.0)}
  \end{equation*}
  where, $M$ is the mass of the central compact object. Rearranging,
  \[
  M \lesssim \frac{c^3\,\delta t}{10G}
  \]
  Thus, upper limit on the mass of central compact core of quasar
  \begin{align*}
  M_{c,\mathrm{max}} &= \frac{c^3\,\delta t}{10G}
  \tag*{(1.0)} \\
  &= \frac{(2.998 \times 10^{8})^3 \times 3600}
          {10 \times 6.674 \times 10^{-11}}~\mathrm{kg}
   = 7.311 \times 10^{7}\,M_\odot \\
  M_{c,\mathrm{max}} &= 7 \times 10^{7}\,M_\odot
  \tag*{(1.0)}
  \end{align*}
\end{description}

\solhead{12.11}{Weighing a galaxy}
% source id: 2011_DA_D2
\ioaafig{2011_DA_D2-s1.pdf}
Solution:

\[ \Delta\lambda = \lambda_1 - \lambda_2 = 20\,\text{\AA} \]
distance on picture ${}\approx 22$ pixels $\pm\,2$ \quad$\Rightarrow$\quad
horizontal scale $20/22 = 0.9091$\,\AA/pixel

Doppler shift from spectrum between upper and lower arm:
$\Delta\lambda = 7$ pixels $\pm\,1 = 6.3636$\,\AA
\[ 2V_{\mathrm{radial}} = c\cdot\Delta\lambda/\lambda \quad
   (c \approx 300\,000\ \mathrm{km/s}) = 290.8\ \mathrm{km/s} \]
Inclination problem: ($i$ -- angle between line of sight and galactic disc)

Size of galaxy image: $a = 11.5 \pm 0.5$\,cm, $b = 6.5 \pm 0.5$\,cm
\quad$\Rightarrow$\quad $i = \arcsin(b/a) \approx 35^\circ$
\[ \cos(i) = V_{\mathrm{radial}}/V_{\mathrm{rotation}}
   \quad\Rightarrow\quad
   V_{\mathrm{rotation}} = V_{\mathrm{radial}}/2\cdot\cos(35^\circ)
   = 216.7\ \mathrm{km/s}
   % sic: as printed; (290.8/2)/cos(35 deg) = 177.5 km/s, not 216.7
\]

Vertical scale (pixel to distance-from-centre conversion)

For 1 pixel: angle $0.82''$ corresponds to
$2\cdot 40\ \mathrm{Mpc}\cdot\mathrm{tg}(0.82''/2) = 159$\,pc in galactic
distance

Mass determination:

From Kepler's law: $M = V^2\cdot R\,/\,G$
\[ R \approx 10\ \mathrm{pixel} = 1590\ \mathrm{pc}\cdot
   3.09\cdot 10^{16}\ \mathrm{m/pc} = 4.9\cdot 10^{19}\ \mathrm{m} \]
\[ M = 3.4\cdot 10^{40}\ \mathrm{kg} \;\Rightarrow\;
   1.7\cdot 10^{10}\,M_\odot \quad
   (M_\odot \approx 2\cdot 10^{30}\ \mathrm{kg}) \]

\solhead{12.12}{Type IA Supernovae}
% source id: 2016_DA_D3
\begin{description}
\item[(D3.1)] Redshifts for the three supernovae are $z_1 = 0.0151$,
  $z_2 = 0.0314$ and $z_3 = 0.0402$. \hfill(1.5 points)

  Filling in the three tables ($\Delta t_{\mathrm{gal}}$, third column; the
  $\Delta t_{\mathrm{s}}$ column is filled in part D3.3): \hfill(3.5 points)

  \begin{center}
  SN2006TD

  \begin{tabular}{rrrr}
    $\Delta t_{\mathrm{obs}}$ (d) & $m_{\mathrm{obs}}$ (mag) &
    $\Delta t_{\mathrm{gal}}$ (d) & $\Delta t_{\mathrm{s}}$ (d) \\
    \hline
    $-15.00$ & 19.41 & $-14.78$ & $-20.00$ \\
    $-10.00$ & 17.48 & $-9.85$  & $-13.34$ \\
    $-5.00$  & 16.12 & $-4.93$  & $-6.67$  \\
    0.00     & 15.74 & 0.00     & 0.00     \\
    5.00     & 16.06 & 4.93     & 6.67     \\
    10.00    & 16.72 & 9.85     & 13.34    \\
    15.00    & 17.53 & 14.78    & 20.00    \\
    20.00    & 18.08 & 19.70    & 26.67    \\
    25.00    & 18.43 & 24.63    & 33.34    \\
    30.00    & 18.64 & 29.56    & 40.01    \\
  \end{tabular}

  \medskip
  SN2006IS

  \begin{tabular}{rrrr}
    $\Delta t_{\mathrm{obs}}$ (d) & $m_{\mathrm{obs}}$ (mag) &
    $\Delta t_{\mathrm{gal}}$ (d) & $\Delta t_{\mathrm{s}}$ (d) \\
    \hline
    $-15.00$ & 18.35 & $-14.54$ & $-14.54$ \\
    $-10.00$ & 17.26 & $-9.70$  & $-9.70$  \\
    $-5.00$  & 16.42 & $-4.85$  & $-4.85$  \\
    0.00     & 16.17 & 0.00     & 0.00     \\
    5.00     & 16.41 & 4.85     & 4.85     \\
    10.00    & 16.82 & 9.70     & 9.70     \\
    15.00    & 17.37 & 14.54    & 14.54    \\
    20.00    & 17.91 & 19.39    & 19.39    \\
    25.00    & 18.39 & 24.24    & 24.24    \\
    30.00    & 18.73 & 29.09    & 29.09    \\
  \end{tabular}

  \medskip
  SN2005LZ

  \begin{tabular}{rrrr}
    $\Delta t_{\mathrm{obs}}$ (d) & $m_{\mathrm{obs}}$ (mag) &
    $\Delta t_{\mathrm{gal}}$ (d) & $\Delta t_{\mathrm{s}}$ (d) \\
    \hline
    $-15.00$ & 20.18 & $-14.42$ & $-17.03$ \\
    $-10.00$ & 18.79 & $-9.61$  & $-11.35$ \\
    $-5.00$  & 17.85 & $-4.81$  & $-5.68$  \\
    0.00     & 17.58 & 0.00     & 0.00     \\
    5.00     & 17.72 & 4.81     & 5.68     \\
    10.00    & 18.24 & 9.61     & 11.35    \\
    15.00    & 18.98 & 14.42    & 17.03    \\
    20.00    & 19.62 & 19.23    & 22.70    \\
    25.00    & 20.16 & 24.03    & 28.38    \\
    30.00    & 20.48 & 28.84    & 34.06    \\
  \end{tabular}
  \end{center}

  The light curves in galaxy frame would appear as follows.
  \hfill(10.0 points)
  \ioaafig{2016_DA_D3-s1.pdf}

\item[(D3.2)] From the graph D3.1, SN2006IS took 22.0\,d to fade by 2
  magnitudes. That is, $\Delta t_0(\mathrm{SN2006IS}) = 22.0$\,d.

  Similarly, $\Delta t_0(\mathrm{SN2006TD}) = 16.4$\,d.

  And $\Delta t_0(\mathrm{SN2005LZ}) = 18.8$\,d.

  \textbf{Acceptable range:} $\pm 1.0$ days \hfill(3.0 points)

  Thus, stretching factors for these two supernovae are
  \begin{align*}
    s_1 &= \frac{22.2}{16.4} = 1.354 \\ % sic: 22.2 as printed, although
                                        % Delta t_0 is quoted as 22.0 d above
    s_3 &= \frac{22.2}{18.8} = 1.181
  \end{align*}
  \hfill(2.0 points)

\item[(D3.3)] Filling the scaled values in the fourth column of the table
  ($\Delta t_{\mathrm{s}}$ in the tables above). \hfill(3.5 points)

  The scaled light curves would appear as follows; the curves should show
  identical profiles. \hfill(10.5 points)
  \ioaafig{2016_DA_D3-s2.pdf}

\item[(D3.4)] To get $a$ and $b$,
  \begin{align*}
    M_{\mathrm{peak},1} &= m_{\mathrm{peak},1} - \mu_1
      = 15.74 - 34.27\ \mathrm{mag} \\
      &= -18.53\ \mathrm{mag} \\
    M_{\mathrm{peak},2} &= m_{\mathrm{peak},2} - \mu_2
      = 16.17 - 35.64\ \mathrm{mag} \\
      &= -19.47\ \mathrm{mag}
  \end{align*}
  \hfill(2.0 points)
  \begin{align*}
    \therefore b &= \frac{s_1 - s_2}{M_{\mathrm{peak},1} - M_{\mathrm{peak},2}}
      = \frac{1.354 - 1}{-18.53 - (-19.47)}\ \mathrm{mag}^{-1}
      = \frac{0.354}{0.94}\ \mathrm{mag}^{-1} \\
    b &= 0.3762\ \mathrm{mag}^{-1}
  \end{align*}
  \hfill(2.0 points)
  \begin{align*}
    a &= s_2 - b\,M_{\mathrm{peak},2} = 1 - 0.3762\times(-19.47)
       = 1 + 7.325 \\
    a &= 8.325
  \end{align*}
  \hfill(2.0 points)

\item[(D3.5)] \mbox{}
  \begin{align*}
    s_3 &= a + b\,M_{\mathrm{peak},3} \\
    \therefore M_{\mathrm{peak},3} &= \frac{s_3 - a}{b}
      = \frac{1.181 - 8.325}{0.3762}\ \mathrm{mag}
      = \frac{-7.144}{0.3762}\ \mathrm{mag} \\
    M_{\mathrm{peak},3} &= -18.99\ \mathrm{mag}
  \end{align*}
  \hfill(2.0 points)

  Distance modulus to SN2005LZ is
  \begin{align*}
    \mu_3 &= m_{\mathrm{peak},3} - M_{\mathrm{peak},3}
      = 17.58 - (-18.99)\ \mathrm{mag} \\
    \mu_3 &= 36.57\ \mathrm{mag}
  \end{align*}
  \hfill(2.0 points)

\item[(D3.6)] Distance to SN2005LZ is
  \begin{align*}
    d_3 &= 10^{\left(\frac{\mu_3}{5}+1\right)}\,\mathrm{pc}
        = 10^{\left(\frac{\mu_3}{5}+1-6\right)}\,\mathrm{Mpc} \\
        &= 10^{\left(\frac{36.57}{5}-5\right)}\,\mathrm{Mpc}
        = 10^{2.314}\,\mathrm{Mpc} \\
        &\simeq 206\ \mathrm{Mpc}
  \end{align*}
  \begin{align*}
    H_0 &= \frac{c z_3}{d_3}
      = \frac{12060}{206}\ \mathrm{km\,s^{-1}\,Mpc^{-1}} \\
    H_0 &= 58.5\ \mathrm{km\,s^{-1}\,Mpc^{-1}}
  \end{align*}
  \hfill(4.0 points)
  \begin{align*}
    T_H &= \frac{1}{H_0}
      = \frac{3.086\times 10^{22}}{58.5\times 10^{3}\times 3.156\times 10^{7}}
      \ \mathrm{yr} \\
    T_H &= 16.7\ \mathrm{Gyr}
  \end{align*}
  \hfill(2.0 points)

  Extra factor of $2/3$ allowed in the value of $T_H$.
\end{description}

\solhead{12.13}{Calibrating the distance ladder to the LMC}
% source id: 2017_DA_D1
\begin{description}
\item[(D1.1)] For Cepheid variable the PL relation is
  \[ \log L = \beta\log P + C \]
  and recall that $F = \dfrac{L}{4\pi d^2}$, so $\log L$ can be written in
  term of absolute magnitude, $M_x$:
  \[ \log L = \log F + 2\log d + \log 4\pi \]
  and
  \[ -\frac{m}{2.5} = \log F - \log F_0 \]
  subtract the 2 equations, we get
  \[ 2.5\log L = -m + 5\log d + C^{*} \]
  And from definition of absolute magnitude,
  \[ M = m + 5 - 5\log_{10} d_{\mathrm{pc}} \]
  We therefore get the relation between $L$ vs.\ $M$,
  \[ 2.5\log L = -M + C'' \]
  Substitute in the PL relation, we get
  \[ M = \beta'\log P + C' \quad\text{or via realising that }
     (\log L \propto M) \]
  where
  \[ \beta' = -2.5\beta \]

  Calculate distances and uncertainties from
  \[ d_{\mathrm{pc}}(\mathrm{parsec}) \approx
     \frac{1\,\mathrm{AU}}{\theta_{\mathrm{parallax}}(\mathrm{arcsec})},
     \qquad
     \Delta d_{\mathrm{pc}} = \frac{\Delta\theta}{\theta}\times
     d_{\mathrm{pc}} \]
  And then calculate absolute magnitude and its uncertainty using
  \[ M_x = m_x - A_x + 5 - 5\log_{10} d_{\mathrm{pc}} \]
  \[ \Delta M_x = \frac{5}{d_{\mathrm{pc}}\ln 10}\times
     \Delta d_{\mathrm{pc}} \]
  $\Delta M_K$ follows from realizing that $\Delta M_K = \Delta M_V$.

  \begin{center}
  \begin{tabular}{lrrrrrrr}
    & $d_{\mathrm{pc}}$ & $\Delta d_{\mathrm{pc}}$ & $\log P$ &
    $M_V$ & $\Delta M_V$ & $M_K$ & $\Delta M_K$ \\
    \hline
    RT Aur      & 416.6 & 32. & 0.572 & $-2.83$ & 0.17 & $-4.19$ & 0.17 \\
    FF Aql      & 355.8 & 22. & 0.650 & $-3.02$ & 0.13 & $-4.37$ & 0.13 \\
    X Sgr       & 333.3 & 20. & 0.846 & $-3.63$ & 0.13 & $-5.12$ & 0.13 \\
    $\zeta$ Gem & 359.7 & 23. & 1.007 & $-3.92$ & 0.14 & $-5.69$ & 0.14 \\
    l Car       & 497.5 & 49. & 1.551 & $-5.27$ & 0.21 & $-7.47$ & 0.21 \\
  \end{tabular}
  \end{center}

  Plot graph $\log P$ vs absolute magnitudes for each band, with error bars,
  and draw a straight line through the data points.
  \ioaafig[0.68]{2017_DA_D1-s1.pdf}
  \ioaafig[0.66]{2017_DA_D1-s2.pdf}

  V-band: slope $= -2.5 \pm 0.2$, with uncertainty between 0.1--0.3.

  K-band: slope $= -3.4 \pm 0.2$, with estimated uncertainty between
  0.1--0.3.

\item[(D1.2)] Calculate the Wesenheit magnitude from the given $V$, $I$ and
  $R_V$ values and then estimate the absolute magnitude for the Wesenheit
  system $M_{W_{VI}}$:

  \begin{center}
  \begin{tabular}{lrr}
    & $W_{VI}$ & $M_{W_{VI}}$ \\
    \hline
    RT Aur      & 3.78 & $-4.31$ \\
    FF Aql      & 3.26 & $-4.49$ \\
    X Sgr       & 2.36 & $-5.25$ \\
    $\zeta$ Gem & 1.88 & $-5.89$ \\
    l Car       & 0.85 & $-7.63$ \\
  \end{tabular}
  \end{center}

  Realising that $\Delta M_{W_{VI}} = \Delta M_V = \Delta M_K$ (or by full
  calculation). Plot $\log P$ vs $M_{W_{VI}}$ with error bars and draw an
  appropriate straight line through the data.
  \ioaafig[0.65]{2017_DA_D1-s3.pdf}

  Estimated slope $= -3.4 \pm 0.2$, with estimated uncertainty of slope
  between 0.1--0.3.

\item[(D1.3)] Distance modulus in each band is given by
  \[ \mu_x = m_x - M_x \]
  And from previous section
  \[ M_x = \beta'_x \log P + C'_x , \]
  where we have estimated slope $\beta$ but not the intercept $C'_x$.

  Therefore, to get the distance in parsec we will need to evaluate
  \[ \mu = m_x - (\beta_x\log P + C'_x) = 5\log_{10} d_{\mathrm{pc}} - 5 \]
  Using the best-estimate slope for each band and any point along the
  best-fit line (or intercept) one can estimate $C'_x$:

  V-band: $C'_V = -1.4$

  $W_{VI}$ band: $C'_{W_{VI}} = -2.4$

  \begin{center}
  \begin{tabular}{lrrrrr}
    Star & $\log P$ & $\mu_V$ & $\mu_{W_{VI}}$ &
    \begin{tabular}{@{}c@{}}V distance\\(pc)\end{tabular} &
    \begin{tabular}{@{}c@{}}$W_{VI}$ distance\\(pc)\end{tabular} \\
    \hline
    HV12199 & 0.42 & 18.53 & 18.39 & 50813 & 47596 \\
    HV12203 & 0.47 & 18.50 & 18.40 & 50223 & 47805 \\
    HV12816 & 0.96 & 18.10 & 18.46 & 41640 & 49221 \\
    HV899   & 1.49 & 18.19 & 18.43 & 43549 & 48660 \\
    HV2257  & 1.60 & 18.25 & 18.36 & 44620 & 47058 \\
    \hline
    & & & Mean & 46170 & 48068 \\
    & & & STD. & 4116  & 865   \\
  \end{tabular}
  \end{center}

  \begin{description}
  \item[a)] Answer NO, the distances estimated from V and Wesenheit are not
    statistically different (they are within one standard deviation of one
    another).
  \item[b)] Answer YES, standard deviation in the estimated distances from
    the 2 bands are significantly different.
  \item[c)] Wesenheit magnitude is better at estimating the distance to LMC.
  \end{description}
\end{description}

\solhead{12.14}{Dust and Young Stars in Star-forming Galaxies}
% source id: 2018_DA_D1
1.1. UV continuum slope ($\beta$) \hfill(23 marks)

\begin{description}
\item[(D1.1.1)] \hfill(6 marks)\\
  The AB magnitude of a typical galaxy is roughly constant, therefore:
  \[ f_\nu = \mathit{const.} \quad\ldots\ldots(1) \]
  Conversion between $f_\nu$ and $f_\lambda$ could be expressed as:
  \[ f_\nu\,\Delta\nu = f_\nu\left(\frac{c}{\lambda-\Delta\lambda}
     - \frac{c}{\lambda}\right)
     = f_\nu\,\frac{c}{\lambda^2}\,\Delta\lambda = f_\lambda\,\Delta\lambda
     \;\Rightarrow\; f_\lambda \propto \lambda^{-2} \quad\ldots\ldots(2) \]
  Therefore, the UV continuum slope should be $\beta = -2$.

\item[(D1.1.2)] \hfill(12 marks)\\
  First of all, we need to convert all the wavelength into rest-frame value,
  using the definition of redshift:
  \[ z = \frac{\lambda_{\mathrm{obs}}}{\lambda_{\mathrm{rest}}} - 1
     \;\Rightarrow\;
     \lambda_0 = \frac{\lambda_{\mathrm{obs}}}{1+z} \quad\ldots\ldots(3) \]
  Thus, the rest-frame wavelengths of the four bands are:

  \begin{center}
  \begin{tabular}{lcccc}
    Band & Y & J & H & K \\
    \hline
    Rest-frame Wavelength (\AA) & 1382 & 1645 & 2171 & 2829 \\
  \end{tabular}
  \end{center}

  The plot is shown as Figure 1.
  \ioaafig{2018_DA_D1-s1.pdf}

\item[(D1.1.3)] \hfill(5 marks)\\
  Here we can express the AB magnitude of CR7 (UV continuum only) using
  $\beta$ and $\lambda_{\mathrm{rest}}$:
  \[ m_{AB} = -2.5\log f_\nu + C = -2.5\log(\lambda^2 f_\lambda) + C
     = -2.5(\beta+2)\log\!\left(\frac{\lambda}{1600\,\text{\AA}}\right)
       + m_{1600} \quad\ldots\ldots(4) \]
  Here $m_{1600}$ is the AB magnitude at rest-frame 1600\,\AA\
  (observed-frame $(1+z)\cdot 1600$\,\AA). After writing this out, we can fit
  the data with a linear function $f(\log\lambda) = a\log\lambda + b$, where
  the slope $a = -2.5(\beta+2)$. Least squares fitting result is $a = 1.7$,
  therefore $\beta = -2.7$.

  The best-fit UV continuum should be plot as a straight line on Figure 1.

  Compared to the galaxy in (D1.1), CR7 shows smaller $\beta$ value,
  indicating its UV continuum is bluer than typical galaxy. Therefore, there
  should be less dust content in CR7, and the correct answer should be [NO].
  \hfill(2 marks)
\end{description}

1.2. UV slope and IRX \hfill(20 marks)

\begin{description}
\item[(D1.2.1)] \hfill(14 marks)\\
  For the galaxy sample in Table 2, the IRX could be calculated as:
  \[ \mathrm{IRX} = \log F_{FIR} - \log F_{1600} \quad\ldots\ldots(5) \]
  Therefore, the final dataset we need for plotting is:

  \begin{center}
  \begin{tabular}{lrr|lrr}
    Galaxy Name & $\beta$ & IRX & Galaxy Name & $\beta$ & IRX \\
    \hline
    NGC 4861     & $-2.46$ & $-0.08$ & NGC 7250 & $-1.33$ & 0.46 \\
    Mrk 153      & $-2.41$ & $-0.55$ & NGC 7714 & $-1.23$ & 0.84 \\
    Tol 1924-416 & $-2.12$ & $-0.12$ & NGC 3049 & $-1.14$ & 0.85 \\
    UGC 9560     & $-2.02$ & $-0.03$ & NGC 3310 & $-1.05$ & 1.01 \\
    NGC 3991     & $-1.91$ & 0.34    & NGC 2782 & $-0.9$  & 1.17 \\
    Mrk 357      & $-1.8$  & 0.21    & NGC 1614 & $-0.76$ & 2.07 \\
    Mrk 36       & $-1.72$ & $-0.26$ & NGC 6052 & $-0.72$ & 1.14 \\
    NGC 4670     & $-1.65$ & 0.17    & NGC 3504 & $-0.56$ & 1.45 \\
    NGC 3125     & $-1.49$ & 0.55    & NGC 4194 & $-0.26$ & 1.63 \\
    UGC 3838     & $-1.41$ & 0.26    & NGC 3256 & 0.16    & 1.88 \\
  \end{tabular}
  \end{center}

  The plot is shown as Figure 2. Best-fit function is:
  \[ \mathrm{IRX} = 0.98\cdot\beta + 1.96 \quad\ldots\ldots(6) \]
  which should be plotted as a straight line on Figure 2.
  \ioaafig{2018_DA_D1-s2.pdf}

\item[(D1.2.2)] \hfill(6 marks)\\
  With the IRX--$\beta$ relation we derived above, we can calculate the
  difference between observed IRX and expected value (inferred by $\beta$) of
  each galaxy:

  \begin{center}
  \begin{tabular}{lrrr|lrrr}
    Name & $\mathrm{IRX}_{\mathrm{obs}}$ & $\mathrm{IRX}_{\mathrm{exp}}$ &
    $\Delta$IRX & Name & $\mathrm{IRX}_{\mathrm{obs}}$ &
    $\mathrm{IRX}_{\mathrm{exp}}$ & $\Delta$IRX \\
    \hline
    NGC 4861     & $-0.08$ & $-0.45$ & 0.37    & NGC 7250 & 0.46 & 0.66 & $-0.20$ \\
    Mrk 153      & $-0.55$ & $-0.40$ & $-0.15$ & NGC 7714 & 0.84 & 0.75 & 0.09 \\
    Tol 1924-416 & $-0.12$ & $-0.12$ & 0.00    & NGC 3049 & 0.85 & 0.84 & 0.01 \\
    UGC 9560     & $-0.03$ & $-0.02$ & $-0.01$ & NGC 3310 & 1.01 & 0.93 & 0.08 \\
    NGC 3991     & 0.34    & 0.09    & 0.25    & NGC 2782 & 1.17 & 1.08 & 0.09 \\
    Mrk 357      & 0.21    & 0.20    & 0.01    & NGC 1614 & 2.07 & 1.22 & 0.85 \\
    Mrk 36       & $-0.26$ & 0.27    & $-0.53$ & NGC 6052 & 1.14 & 1.25 & $-0.11$ \\
    NGC 4670     & 0.17    & 0.34    & $-0.17$ & NGC 3504 & 1.45 & 1.41 & 0.04 \\
    NGC 3125     & 0.55    & 0.50    & 0.05    & NGC 4194 & 1.63 & 1.71 & $-0.08$ \\
    UGC 3838     & 0.26    & 0.58    & $-0.32$ & NGC 3256 & 1.88 & 2.12 & $-0.24$ \\
  \end{tabular}
  \end{center}

  Therefore, the dispersion of IRX is:
  \[ \sigma = \sqrt{\frac{\sum_{i=1}^{N}(\Delta\mathrm{IRX}_i)^2}{N-1}}
     = 0.28\ \mathrm{dex} \quad\ldots\ldots(7) \]
\end{description}

1.3. IRX and $A_{1600}$ \hfill(20 marks)

\begin{description}
\item[(D1.3.1)] \hfill(6 marks)\\
  Based on the two assumptions we made in Question (3), we can simplify
  $F_{FIR}/F_{1600}$ as:
  \[ \frac{F_{FIR}}{F_{1600}} \approx 10^{0.4 A_{1600}} - 1
     \quad\ldots\ldots(8) \]
  \[ \therefore A_{1600} = 2.5\log\left(1 + 10^{\mathrm{IRX}}\right)
     \quad\ldots\ldots(9) \]

\item[(D1.3.2)] \hfill(12 marks)\\
  Now we can calculate $A_{1600}$ of all the galaxies in Table 2 with
  Equation 9:

  \begin{center}
  \begin{tabular}{lrr|lrr}
    Name & $\beta$ & $A_{1600}$ & Name & $\beta$ & $A_{1600}$ \\
    \hline
    NGC 4861     & $-2.46$ & 0.66 & NGC 7250 & $-1.33$ & 1.47 \\
    Mrk 153      & $-2.41$ & 0.27 & NGC 7714 & $-1.23$ & 2.25 \\
    Tol 1924-416 & $-2.12$ & 0.61 & NGC 3049 & $-1.14$ & 2.27 \\
    UGC 9560     & $-2.02$ & 0.72 & NGC 3310 & $-1.05$ & 2.63 \\
    NGC 3991     & $-1.91$ & 1.26 & NGC 2782 & $-0.9$  & 3.00 \\
    Mrk 357      & $-1.8$  & 1.05 & NGC 1614 & $-0.76$ & 5.18 \\
    Mrk 36       & $-1.72$ & 0.48 & NGC 6052 & $-0.72$ & 2.93 \\
    NGC 4670     & $-1.65$ & 0.99 & NGC 3504 & $-0.56$ & 3.66 \\
    NGC 3125     & $-1.49$ & 1.64 & NGC 4194 & $-0.26$ & 4.10 \\
    UGC 3838     & $-1.41$ & 1.13 & NGC 3256 & 0.16    & 4.71 \\
  \end{tabular}
  \end{center}

  Then we can plot the $A_{1600}$--$\beta$ relation as Figure 3. The best-fit
  equation is:
  \[ A_{1600} = 1.92\cdot\beta + 4.62 \quad\ldots\ldots(10) \]
  \ioaafig{2018_DA_D1-s3.pdf}

\item[(D1.3.3)] \hfill(2 marks)\\
  If Equation 10 is correct, then for a dust-free galaxy, the dust
  attenuation $A_{1600}$ should equal to 0. Therefore, we can derive
  $\beta_0 \approx -2.4$, which is the expect UV continuum slope of a
  ``dust-free'' galaxy.

  (Note: The $\beta_0$ result here is conflict with some observations, for
  example, the UV slope of CR7 which we just derived. In fact, the linear
  model here is not necessarily correct and it is just a simplification of
  complex model: we have made two very rough assumption to derive the
  $A_{1600}-\mathrm{IRX}$ relation, and other physical parameters of galaxy,
  e.g.\ metallicity, will also influence the overall $A_{1600}-\beta$
  relation.)
\end{description}

1.4. Extend to high-redshift universe \hfill(12 marks) (discarded)

\begin{description}
\item[(D1.4.1)] \hfill(6 marks)\\
  Since we have already known the upper limit of $F_{FIR}$ of CR7, we only
  need to derive $F_{1600}$ in order to calculate IRX. First of all, we need
  to translate the $m_{1600}$ we got in (D 1.3) to
  $f_{(1+z)1600\text{\AA}}$:
  \[ f_{\lambda_0(1+z)} = f_{\nu_0(1+z)}\cdot
     \frac{c}{[\lambda_0(1+z)]^2}
     = \frac{c}{[\lambda_0(1+z)]^2}\cdot 10^{-0.4 m_{1600}}\cdot
     3631\ \mathrm{Jy} \quad\ldots\ldots(11) \]
  and we will derive
  $f_{\lambda_0(1+z)} = 9.3\times 10^{-13}\,
  \mathrm{J\,s^{-1}\,m^{-3}}$. Therefore:
  \[ F_{1600} = (1+z)\cdot\lambda_0\cdot f_{\lambda_0\cdot(1+z)}
     = 1.1\times 10^{-18}\,\mathrm{J\,s^{-1}\,m^{-2}}
     \quad\ldots\ldots(12) \]
  Thus, the upper limit of CR7's IRX should be:
  \[ \mathrm{IRX}_{CR7} = \log\frac{F_{FIR}}{F_{1600}} = -0.87
     \quad\ldots\ldots(13) \]
  It is an upper limit.

\item[(D1.4.2)] \hfill(6 marks)\\
  Given the UV continuum slop of CR7 ($-2.7$), we can infer its expected IRX
  from the IRX$\,-\,\beta$ relation in local universe:
  \[ \widehat{\mathrm{IRX}}_{CR7} = 0.98\cdot\beta_{CR7} + 1.96 = -0.69
     \quad\ldots\ldots(14) \]
  The upper limit of observed IRX is 0.18\,dex lower than the value predicted
  by local relation. However, we have already known the local relation has a
  dispersion of 0.28\,dex, which is larger than the 0.18\,dex difference.
  This indicates that the current IRX upper limit is not tight enough (in
  other words, current observation is not deep enough) to show any difference
  of CR7 from local galaxies on the diagram of IRX$-\beta$. Therefore, the
  correct answer should be [NO].
\end{description}

\solhead{12.15}{AGN}
% source id: 2020_DA_1
\begin{description}
\item[(a)] Taking multiple reference points, we get that the BLR emission
  lags by about \textbf{20--25 days}. \hfill(1.0 points)

  Answers from 15 to 25 days are acceptable.

\item[(b)] As the AGN is located very far, the time lag can be approximated
  as the time taken by UV emission to reach BLR region. Thus,
  \begin{align*}
    R_{\mathrm{BLR}} = c\Delta t
      &= 3\times 10^{8}\times 20\times 86\,400 \\
      &= 5.2\times 10^{14}\,\mathrm{m} = (0.017\pm 0.004)\,\mathrm{pc}
  \end{align*}
  \hfill(3.0 points)

\item[(c)] As this AGN is 100\,Mpc away from us,
  \begin{align*}
    \theta_{\mathrm{BLR}} &= \frac{0.017}{100e6}\times 206265 \\
      &= (3.5\pm 0.9)\times 10^{-5}\,\mathrm{arcsec}
  \end{align*}
  \hfill(2.0 points)

\item[(d)] The FWHM is approximately $(85\pm 5)$\,\AA\ and the peak is
  approximately at $(4940\pm 5)$\,\AA. \hfill(2.0 points)
  \[ v_\sigma = \frac{\sigma c}{\lambda_{\mathrm{peak}}}
     = \frac{\mathrm{FWHM}\times c}{2.35\,\lambda_{\mathrm{peak}}}
     = \frac{85\times 3\times 10^{5}}{2.35\times 4940} \]
  \hfill(2.0 points)
  \[ v_\sigma = (2200\pm 140)\,\mathrm{km\,s^{-1}} \]
  \hfill(1.0 points)

\item[(e)] \mbox{}
  \begin{align*}
    M_{\mathrm{vir,BH}} &= \frac{v_\sigma^2 R_{\mathrm{BLR}}}{G}
      = \frac{(2.2\times 10^{6})^2\times 5.2\times 10^{14}}
             {6.674\times 10^{-11}} \\
      &= 3.8\times 10^{37}\,\mathrm{kg} \\
    M_{\mathrm{vir,BH}} &= (1.9\pm 0.7)\times 10^{7}\,M_\odot
  \end{align*}
  \hfill(4.0 points)
\end{description}

\solhead{12.16}{Distance to the Large Magellanic Cloud}
% source id: 2023_DA_Q1
\textbf{Distance calculation}

The information from the table can be used to derive individual magnitudes of
both components according to equations:
\begin{align*}
  m_1 &= m_{1+2} + 2.5\log_{10}(1 + L_2/L_1) \\
  m_2 &= m_{1+2} + 2.5\log_{10}(1 + L_1/L_2).
\end{align*}
We will use the third system (OGLE LMC-ECL-18365) as an example to
demonstrate the calculations in detail. The values for this system are as
follows:

Magnitudes: $m_{V,1} = 16.47$\,mag, $m_{K,1} = 14.20$\,mag,
$m_{V,2} = 18.19$\,mag, $m_{K,2} = 16.01$\,mag.

Colours: $(m_V - m_K)_1 = 2.27$\,mag, $(m_V - m_K)_2 = 2.18$\,mag.

The second step is to determine the surface brightness $S_V$ for both
components using the figure. A least square fit to the data in the figure
results in a linear function:
\begin{align*}
  S_V &= 1.346\,\big((m_V - m_K) - 2.407\big) + 5.869\ [\mathrm{mag}],
    \text{ or} \\
  S_V &= 1.346\,(m_V - m_K) + 2.629\ [\mathrm{mag}].
\end{align*}
Uncertainties of the coefficients are $1.346 \pm 0.017$\,mag and
$2.63 \pm 0.04$\,mag.

For this system, this gives $S_{V,1} = 5.69$\,mag and $S_{V,2} = 5.57$\,mag.

However, participants will have two other ways of determining $S_V$.

1) The first way is a graphical way by using a ruler and a pencil in order to
draw the `best-fit' line on the figure. Then $S_V$ follows from an
intersection of the $x = (m_V - m_K)$ vertical line and the `best-fit' line.

2) The second way is to determine the coefficients of the best-fit line
$y = ax + b$ by using coordinates of two points on the figure. The points
should be far from each other.

For example, the coordinates of the second and the penultimate points are:
$(x_1, y_1) = (2.07, 5.41)$ and $(x_2, y_2) = (2.71, 6.28)$. This results in:
\begin{align*}
  a &= \frac{y_2 - y_1}{x_2 - x_1} = 1.36 \\
  b &= \frac{y_1 x_2 - y_2 x_1}{x_2 - x_1} = 2.60
\end{align*}
Using these coefficients we have:
$S_{V,1} = y = 2.27\cdot a + b = 5.69$\,mag and
$S_{V,2} = y = 2.18\cdot a + b = 5.56$\,mag, so within 0.01\,mag of the
`precise' results.

The third step is to calculate angular diameters of components by using the
equation presented in the problem. By modifying the equation defining $S_V$
we obtain:
\[ \theta = 10^{0.2(S_V - m_V)} \]
Subsequently we get: $\theta_1 = 10^{0.2(5.69-16.47)} = 0.00698$\,mas and
$\theta_2 = 10^{0.2(5.56-18.19)} = 0.00298$\,mas.

The fourth step is to calculate the distance to each target. As components
form a physical binary, their distances should be very similar. This is an
independent check of the method and calculations. As the angles under which
we see stellar discs are very small ($\sin\theta \approx \theta$) we can
safely use a linear relation between the angular and physical diameters of an
object. We therefore calculate the distance $D$ as $D = kR/\theta$, where
$\theta$ is expressed in mas, $R$ in solar radii and $k$ is a conversion
factor. The conversion factor results from a fraction
$(2R_\odot/1\,\mathrm{kpc})/(1\,\mathrm{mas}/1\,\mathrm{rad})
= (2R_\odot/1\,\mathrm{AU}) = 9.30\cdot 10^{-3}$.

The distances for the third system are
$D_1 = 9.30\cdot 10^{-3}\cdot 37.30/0.00698 = 49.70$\,kpc and
$D_2 = 9.30\cdot 10^{-3}\cdot 15.94/0.00298 = 49.75$\,kpc.

We then repeat the calculation following the same scheme for the first and
second systems, obtaining for the first system $D_1 = 49.33$\,kpc and
$D_2 = 49.07$\,kpc, and for the second system $D_1 = 49.56$\,kpc and
$D_2 = 49.30$\,kpc. The unweighted mean of all distances is 49.45\,kpc.

\medskip
\textbf{Uncertainties}

\emph{`Statistical' part.}

The standard deviation of the sample is $s = 0.24$\,kpc. The standard error
of the mean is $s/\sqrt{6} = 0.09$\,kpc.

OR:

The mean distances to the three eclipsing binaries are: 49.73\,kpc,
49.20\,kpc and 49.443\,kpc. The standard deviation is $s = 0.22$\,kpc, and
the standard error of the mean is $s/\sqrt{3} = 0.13$\,kpc.

\emph{`Systematic' part.}

All distances are inversely proportional to angular diameters derived from
the $S_V$\,--\,$(m_V - m_K)$ relation. Thus their accuracy is limited by the
precision of the relation. That gives the `irreducible' part of the error:
$0.008\cdot 49.45 = 0.40$\,kpc.

Finally: the distance is $49.45 \pm 0.09 \pm 0.40$\,kpc; the uncertainty is
completely dominated by the precision of the $S_V$\,--\,$(m_V - m_K)$
relation.
